\documentclass[11pt]{article}

\usepackage[T1]{fontenc}
\usepackage[utf8]{inputenc}
\usepackage{lmodern}
\usepackage[margin=1in]{geometry}
\usepackage{amsmath,amssymb,amsthm,mathtools}
\usepackage{enumitem}
\usepackage{booktabs}
\usepackage{array}
\usepackage{longtable}
\usepackage{xcolor}
\usepackage{microtype}
\usepackage{etoolbox}
\usepackage{url}
\usepackage[colorlinks=true,linkcolor=blue!50!black,citecolor=green!35!black,urlcolor=blue!60!black]{hyperref}
\usepackage[capitalize,nameinlink]{cleveref}

% ---------------------------------------------------------------------------
% Theorem environments
% ---------------------------------------------------------------------------
\theoremstyle{plain}
\newtheorem{theorem}{Theorem}[section]
\newtheorem{lemma}[theorem]{Lemma}
\newtheorem{corollary}[theorem]{Corollary}
\newtheorem{proposition}[theorem]{Proposition}
\newtheorem{conjecture}[theorem]{Conjecture}
\newtheorem{question}[theorem]{Question}
\newtheorem{problem}[theorem]{Problem}
\theoremstyle{definition}
\newtheorem{definition}[theorem]{Definition}
\newtheorem{remark}[theorem]{Remark}
\newtheorem{example}[theorem]{Example}

% The Lean proof is the proof; the text gives a faithful outline of it.
\newenvironment{outline}[1][Outline of the formal proof]{\begin{proof}[#1]}{\end{proof}}
% Results that are not formalized get complete proofs, labelled as such.
\newenvironment{informalproof}{\begin{proof}[Proof (not formalized)]}{\end{proof}}

% ---------------------------------------------------------------------------
% Lean references
%   \lean{name}            a declaration, relative to the namespace HalfRate (declarations of
%                          earlier papers of the series are written with their full names)
%   \formal{file}{names}   the formal counterpart of a numbered result: a file relative to
%                          lean/HalfRate/ and a comma-separated list of declarations
%   \partformal{file}{names}{gap}
%                          a numbered result whose formal counterpart covers all but the
%                          part described in {gap}; the gap is proved or cited in the text
%   \notformalized         a numbered result that has no formal counterpart
%   \cited{where}          a result of an earlier paper used without proof: where it is proved,
%                          and whether it is Lean-verified there
% ---------------------------------------------------------------------------
% Lean names are typeset like paths, with line breaks allowed after "_", "." and "/".
\DeclareUrlCommand\leanname{\def\UrlFont{\ttfamily\upshape}}
\renewcommand{\UrlBreaks}{\do\_\do\.\do\/\do\-}
\newcommand{\lean}[1]{\leanname{#1}}
\newcommand{\leanfile}[1]{\leanname{#1}}
\newcommand{\leanlist}[1]{\begingroup\def\leansep{}%
  \renewcommand*{\do}[1]{\leansep\leanname{##1}\def\leansep{, }}\docsvlist{#1}\endgroup}
\newcommand{\formal}[2]{\par\noindent{\small\upshape\raggedright\textsf{Lean:} \leanlist{#2}\unskip
  \nobreak\hfil\penalty50\hskip1em\hbox{}\nobreak\hfill[\leanfile{#1}]\parfillskip=0pt\par}%
  \smallskip\noindent\ignorespaces}
\newcommand{\partformal}[3]{\par\noindent{\small\upshape\raggedright\textsf{Lean (partial):} \leanlist{#2}\unskip
  \nobreak\hfil\penalty50\hskip1em\hbox{}\nobreak\hfill[\leanfile{#1}]\parfillskip=0pt\par}%
  \par\noindent{\small\upshape\textsf{Not formalized:} #3}\par\smallskip\noindent\ignorespaces}
\newcommand{\notformalized}{\par\noindent{\small\upshape\textsf{Not formalized;} a complete proof is given
  below.}\par\smallskip\noindent\ignorespaces}
\newcommand{\cited}[1]{\par\noindent{\small\upshape\textsf{Cited:} #1}\par\smallskip\noindent\ignorespaces}

% ---------------------------------------------------------------------------
% Notation
% ---------------------------------------------------------------------------
\newcommand{\N}{\mathbb{N}}
\newcommand{\Z}{\mathbb{Z}}
\newcommand{\Q}{\mathbb{Q}}
\newcommand{\R}{\mathbb{R}}
\newcommand{\C}{\mathbb{C}}
\newcommand{\F}{\mathbb{F}}
\newcommand{\T}{\mathsf{T}}
\newcommand{\clog}[1]{\lceil \log_2 #1 \rceil}
\newcommand{\abs}[1]{\lvert #1 \rvert}
\newcommand{\card}[1]{\lvert #1 \rvert}
\newcommand{\eps}{\varepsilon}
\newcommand{\prev}{\mathrm{prev}}
\newcommand{\cost}{\mathrm{cost}}
\newcommand{\gn}{\mathrm{gn}}
\newcommand{\Sh}{\mathrm{Sh}}
\newcommand{\cl}{\mathrm{cl}}
\newcommand{\Aut}{\mathrm{Aut}}
\newcommand{\minrk}{\mathrm{minrk}}
\newcommand{\MAIS}{\mathrm{MAIS}}
\newcommand{\Crux}{\mathrm{Crux}}
\newcommand{\TCP}{\mathrm{TCP}}
\DeclareMathOperator{\rank}{rank}
\DeclareMathOperator{\dist}{dist}
\DeclareMathOperator{\poly}{poly}

% ---------------------------------------------------------------------------
% Additions for this paper. Lean declarations are relative to the namespace HalfRate, and the
% files named in \formal and \partformal lines are relative to lean/HalfRate/.
% ---------------------------------------------------------------------------
\newcommand{\Cay}{\operatorname{Cay}}
\newcommand{\Andr}{\operatorname{And}}
\newcommand{\nustar}{\nu^{*}}
\newcommand{\Bdir}{B_{\mathrm{dir}}}
\newcommand{\Tint}{T_{\mathrm{int}}}
\newcommand{\Ex}{\mathrm{X}_{15}}


\title{Nested Elimination and Half-Rate Bounds\\ for Storage Codes}
\author{Hans Brende \and Claude Opus 5.5 (Anthropic)}
\date{October 2026}

\begin{document}

\maketitle

\begin{abstract}
Can a deterministic multitape Turing machine transpose a $d\times d$ matrix of $\lceil\log_2d\rceil$-bit entries in time
linear in the input length? Earlier papers of this series reduce the impossibility of oblivious linear-time transposers to
depth reduction and an all-code question about storage codes: on bounded-degree recovery relations without short paired
paths, can a storage code beat rate~$\frac12$? We prove the exact half bound $2h(V)\le\sum_vh(\{v\})$ for every valid
polymatroid, and hence $\card C^2\le\card\alpha^{\card V}$ for every storage code, linear or not, over every alphabet, on
large classes of recovery relations, most of them defined by the absence of transitive triangles. The engine is a nested
elimination lemma (Lemma~E), which extends the inequality form of an exact partition identity of the series' circular
memo by outside parents with nested charges, and a decomposition theorem (Theorem~D) that composes certified pieces. Its
instances include triangle-free graphs of maximum degree three (Theorem~N3); loopless relations without transitive
triangles whose underlying graph has maximum degree three (Theorem~U3), or whose vertices of in-degree at least three
induce an acyclic sub-relation, with no degree bound (Theorem~T); loopless relations partitioned into blocks with private
outside parents, each block chordless subcubic or a chorded digon path, without transitive triangles, or internally
certified (Theorem~PC); and every Cayley relation $\mathrm{Cay}(G,\{s,s^{-1},t\})$ without transitive triangles, for every
finite group, when $\{s,s^{-1},t\}$ has three elements and does not contain~$1$ (Corollary~CC). On triangle-free graphs of
maximum degree three the guessing number equals the fractional matching number, and the broadcast rate of index coding is
the number of vertices minus it. We also show where the method stops: in a $15$-vertex relation with in- and out-degrees
at most three the unique minimum nonempty core admits no certificate, some chorded cycles admit no valid sweep, the
Petersen graph admits no global certificate, and on the folded cubes, whose odd girth equals their degree, the half bound
fails, so that a half bound assuming only the absence of short paired paths up to a fixed length must restrict the
degree. Apart from three cited results of the fourth paper, which are Lean-verified there, every numbered theorem,
proposition, lemma, corollary and example is formalized in Lean~4, as are the certificates printed in the remarks. Not
formalized are the remaining claims of the numerical remarks (that the $15$-vertex relation is strongly connected and not
$2$-dicolourable and that its minimum nonempty core is not a piece, that twelve chorded cycles have no valid sweep, and
that the Petersen graph has no global certificate), the cited folded-cube formula, which the fourth paper verifies in
Lean for the Clebsch graph, a cited theorem on odd cycles, and a heuristic remark. These are unconditional storage-code theorems: they constrain
constructions and locate the remaining work for an impossibility proof, but they do not decide whether linear-time
transposition is possible.

\end{abstract}

\begin{quote}\small
\textbf{Contributions.} Hans Brende chose the problem, fixed the verification standard, and directed the research.
Claude Opus 5.5, a large language model developed by Anthropic, developed the mathematics, wrote the Lean
formalization, and wrote this paper under Hans Brende's direction. The engine of the paper, the nested elimination
lemma, generalizes the exact partition identity of the circular memo by Hans Brende and GPT-6 Astra
(OpenAI)~\cite[Theorem~2.1]{CircularMemo2026}, and the treatment of graphs of circular chromatic number below three
reuses that memo's Andrásfai elimination order.
\end{quote}

\tableofcontents
\clearpage

\section{Introduction}\label{sec:intro}

\subsection{The problem}

Can a deterministic multitape Turing machine transpose a $d\times d$ matrix with $b(d)=\max(1,\clog d)$-bit entries,
written row by row, in time linear in the input length $L=d^2b(d)$? This is the statement~(T) of the first paper of this
series~\cite{FirstPaper2026}. By a reduction of Harvey and van der Hoeven~\cite[Cor.~6.9]{HvdH2025}, a linear-time
multiplication machine would give a linear-time transposer, so a proof of $\neg$(T) would give the first superlinear lower
bound for off-line integer multiplication on multitape Turing machines. The question is open, and this paper does not
decide it.

For oblivious machines the series has reduced $\neg$(T) to two statements. The fourth
paper~\cite[Theorem~4.10 and Corollary~4.11]{FourthPaper2026} shows that depth reduction of the time--space dependency
graph by a cheap genie, together with an index-coding lower bound on the side information of the remaining bounded-depth
cones (the \emph{crux}), rules out oblivious linear-time transposers. \emph{Directed Recovery Structures in Tape
Computations}~\cite{DirectedRecovery2026} shows that the cone relations of linear-time schedules have few coordinates
whose input and demand ports are close in the paired incidence graph, and that an affirmative answer to one machine-free
question would supply the crux on those cone relations~\cite[Question~9.1 and Proposition~9.2]{DirectedRecovery2026}: do
storage codes on bounded-degree recovery relations with this paired-distance separation have rate at most
$\frac12+o(1)$, for every alphabet?

This is a storage-code paper, like the series paper \emph{Covers and Local Generation in Linear Storage
Codes}~\cite{Covers2026}. We prove the exact half bound for all codes on large classes of recovery relations without short
paired paths, and show where the method stops. A recovery relation assigns to each coordinate $v$ the set $N(v)$ of
coordinates from which it must be recoverable. For a polymatroid $h$ valid for $N$, that is, with
$h(N(v)\cup\{v\})=h(N(v))$ for every $v$, the bound reads
\[
  2h(V)\;\le\;\sum_{v\in V}h(\{v\}).\tag{$\star$}
\]
Applied to the entropy of a uniformly random codeword, \eqref{eq:star} gives $\card C^2\le\card\alpha^{\card V}$ for
every storage code $C\subseteq\alpha^V$, linear or not, over every alphabet~\cite{FourthPaper2026}: no code beats the
pairing. Our hypotheses are exact rather than asymptotic. The central one is that the relation has no \emph{transitive
triangle} $x\to y\to z$, $x\to z$, where $u\to v$ means $u\in N(v)$; for loopless relations this says exactly that every
coordinate has paired distance at least five (\cref{lem:tt-paired}).

\subsection{Results}

\paragraph{The engine (\cref{sec:engine}).} \emph{Lemma~E} (nested elimination) eliminates a piece $H$ one vertex at a
time. Each vertex must be recovered from the vertices already eliminated, from one record other than its own, and possibly
from one vertex outside the piece; steps that use the same outside vertex must be nested. Then
$2h(H)\le\sum_ih(P_i)+\sum_{w\notin H}I(w;H)$ (\cref{lem:E}). Without outside vertices this is the inequality form of the
exact partition identity of the circular memo~\cite[Theorem~2.1]{CircularMemo2026}; the nested charging is new.
\emph{Theorem~D} (decomposition) contracts certified pieces one after another, each in the relation left by the earlier
ones, and proves \eqref{eq:star} whenever the vertex set has such a decomposition (\cref{thm:D}). Relative $2$-dicolourings
(Lemma~$\Bdir$, \cref{lem:Bdir}) and global certificates, which pull back along homomorphisms (\cref{lem:pullback}), are
the basic pieces.

\paragraph{The main theorem.} Every relation in the following classes has a decomposition, or a global certificate, and
hence satisfies \eqref{eq:star}; every storage code on it satisfies $\card C^2\le\card\alpha^{\card V}$.
\begin{enumerate}[label=(\alph*),leftmargin=2.5em]
\item $2$-dicolourable relations, in particular bipartite graphs (\cref{prop:twoAcyclic});
\item graphs of circular chromatic number below three, at any degree; this is the circular memo's
theorem~\cite[Corollary~3.2]{CircularMemo2026}, re-derived here as a direct instance (\cref{thm:circular});
\item triangle-free graphs of maximum degree three, and every relation whose arcs join adjacent vertices of such a graph
(\emph{Theorem~N3}, \cref{thm:N3});
\item loopless TT-free relations whose underlying graph has maximum degree three, directed triangles included
(\emph{Theorem~U3}, \cref{thm:U3});
\item loopless TT-free relations whose vertices with at least three parents induce an acyclic sub-relation, with no bound
on in- or out-degrees (\emph{Theorem~T}, \cref{thm:T}), and the version in which only parents inside a vertex's strongly
connected component count (\emph{Corollary~$\Tint$}, \cref{cor:Tint});
\item loopless relations partitioned into blocks with private outside parents, each block $K$ chordless subcubic or a
chorded digon path, with $N[K]$ TT-free, or internally certified (\emph{Theorem~PC}, \cref{thm:PC});
\item every TT-free Cayley relation $\Cay(G,\{s,s^{-1},t\})$, with arcs $g\to gx$, for every finite group $G$, when
$\{s,s^{-1},t\}$ has three elements and does not contain~$1$ (\emph{Corollary~CC}, \cref{cor:CC}).
\end{enumerate}
Corollary~$\Tint$ strictly contains the TT-free case of the two-input theorem of \emph{Directed Recovery
Structures}~\cite[Theorem~5.1]{DirectedRecovery2026}, and Theorem~U3 is the all-code version, on subcubic underlying
graphs, of the fourth paper's scalar linear Theorem~S~\cite[Theorem~6.13]{FourthPaper2026}. The certificates of
Theorem~PC come from \emph{Lemma~S} for chorded paths (\cref{lem:S}) and from sweeps of chorded circulant cycles
(\emph{Proposition~CC}, \cref{prop:CC}). The relation $\Cay(\Z_9,\{1,-1,3\})$ lies in none of the classes (a)--(e); it is covered
by~(g), and by~(f) with a single block.

\paragraph{Exact values (\cref{sec:consequences}).} On a triangle-free graph $G$ of maximum degree three, every storage code
satisfies $2\log\card C\le(\card V+\card S-\card{I(S)})\log\card\alpha$ for every vertex set $S$, where $I(S)$ is the set
of isolated vertices of $G-S$ (\cref{thm:barrier}). With the fractional Tutte--Berge formula (\cref{prop:fracTB}), the
guessing number equals the fractional matching number $\nustar(G)$, attained over every alphabet of square size
(\cref{thm:gn,cor:gn}). This proves, for these graphs and with the asymptotic guessing number taken either as the
supremum or as the limit over alphabet sizes, Christofides and Markström's conjecture that fractional clique covers are
optimal guessing strategies~\cite{ChristofidesMarkstrom2011,Dang2015}, which is false in
general~\cite{Baber2017,CameronDangRiis2016}; it was known before for perfect graphs, bipartite graphs in particular, and odd
cycles~\cite{ChristofidesMarkstrom2011}, and for prisms and certain chorded cycles~\cite{MMV2019}. It also proves Dang's
conjecture for planar triangle-free graphs at maximum degree three~\cite[Conjecture~15]{Dang2015}. Dually, the broadcast
rate of index coding is $\card V-\nustar(G)$ (\cref{thm:index,cor:broadcast}). Relations that become one of our classes
after deleting a set $Z$ have storage codes of rate at most
$(\card V+\card Z)/(2\card V)$ (\cref{cor:paid}).

\paragraph{Limits (\cref{sec:limits}).} The natural rules for choosing pieces fail: $\Cay(\Z_9,\{1,-1,3\})$ has no proper
piece with private outside parents (\cref{prop:no-private}); a $15$-vertex relation with in- and out-degrees at most three
has a unique minimum nonempty core, which is not a piece, although the relation has a global certificate
(\cref{rem:X15}); some chorded cycles have no valid sweep (\cref{rem:sweeps}); and the Petersen graph has no global
certificate, although Theorem~N3 gives it a decomposition (\cref{rem:petersen}). The half bound fails on the folded
cubes, whose odd girth equals their degree (\cref{rem:foldcube}, from a rank formula of Barg, Schwartz and
Yohananov~\cite{BSY2024}); so an exact half bound that assumes only $b_E(r)=0$, for a fixed odd $r\ge3$, cannot
hold at maximum degree $r+2$. For $r=3$ it fails already at maximum degree four: a $4$-regular triangle-free graph on
thirteen vertices, of odd girth five, carries a linear storage code over~$\F_4$ with $\card C=4^7$ (\cref{ex:deg4}), so
Theorem~N3 is sharp in the degree.

\paragraph{What this means for (T) (\cref{sec:construction}).} These are unconditional storage-code theorems; they do not
decide~(T). They constrain constructions: in the cone relations of a correct oblivious transposer left by
a genie of cost $\eps L$, no induced sub-relation on more than $2\eps L$ coordinates may lie in any of the classes above.
They also locate the remaining work for an impossibility proof: \cref{conj:P}, that every loopless TT-free relation of in-degree at most three has a decomposition,
would give the degree-three case of the paired-distance question (\cref{prop:dn3-transfer}).

\subsection{How to read this paper: the formal proofs}

Apart from three results of the fourth paper that we cite (\cref{thm:codes,lem:closure,lem:fewer-arcs}), which are
Lean-verified there, every numbered theorem, proposition, lemma, corollary and example of this paper is formalized in
Lean, including the guessing number and the broadcast rate (\cref{cor:gn,cor:broadcast}), the conditional transfer of
\cref{prop:dn3-transfer}, \cref{prop:no-private}, \cref{ex:D9} and \cref{ex:deg4}; the finite checks are evaluated by the
Lean kernel.
Not formalized are the remaining claims of the numerical remarks of \cref{sec:limits}: that the $15$-vertex relation
$\Ex$ is strongly connected and not $2$-dicolourable and that its minimum nonempty core is not a piece; that the twelve
chorded cycles of \cref{rem:sweeps} have no valid sweep; and that the Petersen graph has no global certificate. They are
reproducible with the scripts in \leanfile{tools/}. Not formalized either are the cited rank formula for folded cubes,
whose case of the Clebsch graph is Lean-verified in the fourth paper, the cited theorem of Atkins, Rombach and Skerman
after \cref{cor:gn}, and the heuristic \cref{rem:degree4}. The certificates printed in the remarks are Lean-verified.
None of the unformalized claims is used in a proof. Each numbered result states its status as described in \cref{sec:prelim-lean}.

\subsection{Related work}

\paragraph{Guessing numbers and graph entropy.} Guessing games and guessing numbers were introduced by
Riis~\cite{Riis2007}. Christofides and Markström~\cite{ChristofidesMarkstrom2011} studied undirected graphs, introduced
the fractional clique cover strategy and conjectured it to be optimal; the conjecture was disproved by Baber,
Christofides, Dang, Vaughan and Riis~\cite{Baber2017}, with triangle-free counterexamples by Cameron, Dang and
Riis~\cite{CameronDangRiis2016}. Dang's thesis~\cite{Dang2015} computes guessing numbers of several families of planar
triangle-free graphs; those of maximum degree three are special cases of Theorem~N3. Gadouleau and
Riis~\cite{GadouleauRiis2011} relate the guessing number of a digraph to the independence number of its guessing graph and
construct digraphs with large guessing numbers. The bounds there come from Shannon and non-Shannon inequalities or from
explicit strategies. Mazumdar, McGregor and Vorotnikova~\cite{MMV2019} bound storage capacity by covering a graph with
gadgets of known capacity, which already gives combinatorial certificates for Shannon-type upper bounds on some families.
Our certificates are of a different kind: elimination orders whose steps may use one outside vertex, with nested charges
(Lemma~E), composed by contraction rather than by covering (Theorem~D).

\paragraph{Storage codes and index coding.} Storage codes on graphs and their duality with index coding were studied by
Mazumdar~\cite{Mazumdar2014}, by Shanmugam and Dimakis~\cite{ShanmugamDimakis2014}, and by Arbabjolfaei and
Kim~\cite{ArbabjolfaeiKim2015}. Index coding was introduced by Birk and Kol~\cite{BirkKol2006} and formalized, with minrank
and the acyclic-subgraph lower bound, by Bar-Yossef, Birk, Jayram and Kol~\cite{BBJK2011}; Blasiak, Kleinberg and
Lubetzky~\cite{BKL2013} bound and approximate the broadcast rate by linear programming. Mazumdar, McGregor and Vorotnikova
bounded the storage capacity above by vertex covers and below by fractional matchings~\cite{MMV2019}; the lower bound is
Christofides and Markström's fractional clique cover strategy~\cite{ChristofidesMarkstrom2011} in storage language, and on
triangle-free subcubic graphs it is exact (\cref{thm:gn}). Without a degree bound, triangle-freeness does not force
rate~$\frac12$: Barg and Zémor~\cite{BargZemor2022} and Barg, Schwartz and Yohananov~\cite{BSY2024} construct storage codes
on triangle-free graphs of growing degree with rates approaching $\frac34$ and~$1$, and an earlier triangle-free family of
rate approaching~$1$ is due to Golovnev and Haviv~\cite{GolovnevHaviv2022}, as noted in~\cite{BSY2024}. The
bounded-degree examples among these have degree at least five, such as the Clebsch graph (\cref{rem:foldcube}); we are
not aware of an earlier triangle-free example of maximum degree four (\cref{ex:deg4}), and at degree three there is
none (Theorem~N3).

\paragraph{The series.} The fourth paper~\cite{FourthPaper2026} proved that rate bounds pull back along homomorphisms
(Theorem~H), the half bound for scalar linear codes on triangle-free subcubic side information with its matching form
(Theorem~S), and that the Clebsch graph has guessing number ten. The circular memo~\cite{CircularMemo2026} proved the
exact partition identity on which Lemma~E builds, and the half bound for circular chromatic number below three. The
subcubic memo~\cite{SubcubicMemo2026} proved half bounds for reversible and low-rank codes on subcubic supports and, for
all codes, on triangle-free subcubic graphs covered by four-cycles; Theorem~N3 contains the latter. \emph{Directed
Recovery Structures}~\cite{DirectedRecovery2026} proved the paired-distance bound for cone relations, the two-input theorem
and the hyperfinite half bound, and posed the all-code paired-distance question that motivates this paper. The covers
paper~\cite{Covers2026} studies cover structure and local generation of linear storage codes. The companion
paper~\cite{EighthPaper2026}, published together with this one, is about the depth side: its Theorem~F proves depth
reduction for head walks of bounded nesting load, so the crux alone, or an affirmative answer to Question~9.1 of
\emph{Directed Recovery Structures}, excludes oblivious linear-time transposers whose work heads enter every cell boundedly
often.

To our knowledge, Lemma~E, Theorems~N3, U3, T and~PC, and Corollary~CC are new. Special cases were known: the guessing
numbers of perfect graphs and odd cycles~\cite{ChristofidesMarkstrom2011}; the storage capacity of prisms, of certain
chorded cycles and of odd cycles~\cite{MMV2019}; the broadcast rates of cycles and of the cubic circulants, and numerically
that of the Petersen graph~\cite{BKL2013}; the families of planar triangle-free graphs in Dang's thesis~\cite{Dang2015};
the half bound on square-covered subcubic graphs~\cite{SubcubicMemo2026}; and the half bound for circular chromatic number
below three, which is the circular memo's theorem~\cite{CircularMemo2026}.

\subsection{Organization}

\Cref{sec:prelim} fixes notation. \Cref{sec:engine} proves Lemma~E and Theorem~D and the tools for building
certificates. \Cref{sec:undirected} treats $2$-dicolourable relations and symmetric side information, \cref{sec:directed}
the directed classes, and \cref{sec:consequences} the barrier bound, guessing numbers and index codes. \Cref{sec:limits}
shows where the method stops, \cref{sec:construction} draws the consequences for constructions, and \cref{sec:open} lists
open problems. The appendices index the Lean declarations and explain how to build and check the library.

\section{Preliminaries}\label{sec:prelim}

\subsection{The statement (T)}\label{sec:prelim-T}

We use the machine model of the first paper of this series~\cite{FirstPaper2026}: a deterministic Turing machine with
a read-only input tape, a constant number of work tapes over a finite alphabet, and a write-only output tape. For
$d\ge1$ let $b(d)=\max(1,\clog d)$. The input is a $d\times d$ matrix with $b(d)$-bit entries written row by row, of
length $L=d^2b(d)$, and the machine must write its transpose in the same encoding. Statement~(T) says that some
machine does this within $c(L+1)$ steps, for a constant $c$ and all $d$. Its formal version,
\lean{Transpose.LinearTimeTransposable}, is fixed in a frozen file of the first paper's library. A machine is
\emph{oblivious} if its head positions, output times and halting time depend only on the input length.

The Lean library of this paper does not model machines. Its theorems are about recovery relations, polymatroids,
storage codes and index codes on finite sets. The connection to~(T) goes through the fourth paper of the
series~\cite{FourthPaper2026} and through \emph{Directed Recovery Structures in Tape
Computations}~\cite{DirectedRecovery2026}; we recall it in \cref{sec:construction}.

\subsection{Recovery relations}\label{sec:prelim-relations}

A \emph{recovery relation} on a finite set $V$ is a map $N\colon V\to2^V$. We write $u\to v$, an \emph{arc}, when
$u\in N(v)$: then $u$ is a \emph{parent} (in-neighbour) of $v$ and $v$ a \emph{child} of $u$. The \emph{in-degree}
of $v$ is $\card{N(v)}$. The relation is \emph{loopless} if $v\notin N(v)$ for every $v$. For $K\subseteq V$, the
\emph{induced relation} $N[K]$ on $K$ is $v\mapsto N(v)\cap K$, and $N-Z=N[V\setminus Z]$. For a simple graph $G$ on
$V$, its \emph{neighbourhood relation} $N_G$ maps $v$ to its set of neighbours (\lean{MotionDepthCodes.nbhd}); every
arc of $N_G$ is reciprocated.

A \emph{digon} is a pair of arcs $u\to v$ and $v\to u$; an arc that is not part of a digon is \emph{one-way}. The
\emph{digon graph} $G_2(N)$ is the simple graph whose edges are the digons. The \emph{underlying graph} $U(N)$
(\lean{und}) joins $u\ne v$ whenever $u\to v$ or $v\to u$. A \emph{transitive triangle} is a triple of distinct
vertices $x,y,z$ with $x\to y$, $y\to z$ and $x\to z$; we call $y$ its \emph{middle vertex}. A relation is
\emph{TT-free} if it has no transitive triangle (\lean{TTFree}); \lean{TTFreeOn} is the same condition for triples
inside a set. In a loopless TT-free relation the parents of a vertex are pairwise non-adjacent, the digon graph is
triangle-free, and every triangle of the underlying graph is a \emph{directed $3$-cycle} $x\to y\to z\to x$ without
reverse arcs (\cref{lem:dtri} below); bidirected triangles are excluded.

A set $A\subseteq V$ is \emph{acyclic} for $N$ if every nonempty $T\subseteq A$ contains a vertex with no parent in
$T$ (\lean{IsAcyclicIn}); equivalently, $N[A]$ has no loops and no directed cycles. A relation is
\emph{$2$-dicolourable} if $V$ is the union of two disjoint acyclic sets.

\paragraph{Paired incidence.} \emph{Directed Recovery Structures}~\cite[Definition~2.3]{DirectedRecovery2026} attaches
to $N$ its incidence graph $E(N)$: a bipartite graph with vertices $L_p$ and $R_p$ for $p\in V$ and an edge
$L_pR_q$ whenever $p\in N(q)$. It puts $b_E(r)=\card{\{p\in V:\dist_E(L_p,R_p)\le r\}}$, and a sequence of relations
satisfies the \emph{paired-distance condition} if $b_E(r)=o(\card V)$ for every fixed $r$. Distances between $L_p$
and $R_p$ are odd. Our hypothesis ``TT-free'' is the exact form of this condition at radius three.

\begin{lemma}[Transitive triangles are paired paths of length three]\label{lem:tt-paired}
\notformalized
Let $N$ be a recovery relation and $p\in V$. Then $\dist_E(L_p,R_p)=1$ if and only if $p\in N(p)$. If $N$ is
loopless, then $\dist_E(L_p,R_p)=3$ if and only if $p$ is the middle vertex of a transitive triangle. In particular,
a loopless relation is TT-free if and only if $b_E(3)=0$, that is, if and only if $\dist_E(L_p,R_p)\ge5$ for every
$p$.
\end{lemma}

\begin{informalproof}
The edge $L_pR_p$ exists exactly when $p\in N(p)$. Let $N$ be loopless. A path $L_p,R_z,L_x,R_p$ of length three
means $p\in N(z)$, $x\in N(z)$ and $x\in N(p)$, that is, $p\to z$, $x\to z$ and $x\to p$. Its vertices are distinct,
so $z\ne p$ and $x\ne p$, and $x\ne z$ because $x\in N(z)$ and $N$ is loopless. Hence $x\to p\to z$ with $x\to z$ is
a transitive triangle with middle vertex $p$. Conversely, every transitive triangle $x\to p\to z$, $x\to z$ gives
this path. Since the distance is odd and not $1$, it is at least $5$ when no such path exists.
\end{informalproof}

\subsection{Polymatroids and storage codes}\label{sec:prelim-codes}

We use the definitions of the fourth paper~\cite[Definitions~2.1 and~2.2]{FourthPaper2026}. A \emph{polymatroid} on
$V$ is a function $h\colon2^V\to\R$ with $h(\emptyset)=0$ that is monotone and submodular. It is \emph{valid} for
$N$ if $h(N(v)\cup\{v\})=h(N(v))$ for every $v\in V$. We write $h(v)$ for $h(\{v\})$,
\[
  h(A\mid B)=h(A\cup B)-h(B),\qquad I(A;B\mid S)=h(A\cup S)+h(B\cup S)-h(A\cup B\cup S)-h(S)\ge0,
\]
and $I(w;X)=h(w)+h(X)-h(X\cup\{w\})$ (\lean{mi}). The subject of this paper is the \emph{weighted half bound}
\begin{equation}\label{eq:star}
  2h(V)\;\le\;\sum_{v\in V}h(v)\qquad\text{for every polymatroid $h$ valid for $N$.}\tag{$\star$}
\end{equation}
In the terms of the fourth paper, \eqref{eq:star} says that $N$ satisfies the rate bound $\tfrac12$; the
formal statement is \lean{MotionDepthCodes.RateBound}~$N\,(1/2)$. A \emph{storage code} for $N$ over a finite alphabet $\alpha$ is a set
$C\subseteq\alpha^V$ such that two words of $C$ that agree on $N(v)$ agree at $v$, for every $v$
(\lean{MotionDepthCodes.IsStorageCode}). The \emph{guessing number} $\gn(N,q)$ is the largest value of
$\log_q\card C$ over storage codes with $\card\alpha=q$, and $\gn(N)=\sup_q\gn(N,q)$ is the asymptotic guessing
number~\cite{Riis2007,ChristofidesMarkstrom2011}.

\begin{theorem}[Rate bounds bound codes]\label{thm:codes}
\cited{the fourth paper~\cite[Theorem~6.3 and Corollary~6.9]{FourthPaper2026}, Lean-verified there
(\lean{MotionDepthCodes.card_sq_le_of_rateBound_half}).}
If $N$ satisfies \eqref{eq:star}, then every storage code $C\subseteq\alpha^V$ for $N$ over a nonempty finite
alphabet satisfies $\card C^2\le\card\alpha^{\card V}$.
\end{theorem}

The proof applies \eqref{eq:star} to the entropy of a uniformly random codeword, which is a valid
polymatroid. The same argument gives $\card C^2\le\prod_v\card{\alpha_v}$ for codes with coordinate alphabets
$\alpha_v$~\cite[Corollary~3.4]{CircularMemo2026}; we do not use this.

\begin{lemma}[Closure]\label{lem:closure}
\cited{the fourth paper~\cite[Lemma~6.4]{FourthPaper2026}, Lean-verified there
(\lean{MotionDepthCodes.IsPolymatroid.eq_of_closure}).}
Let $h$ be a polymatroid valid for $N$, and let $S\subseteq T$ be such that $N(v)\subseteq S$ for every
$v\in T\setminus S$. Then $h(T)=h(S)$.
\end{lemma}

\begin{lemma}[Fewer arcs]\label{lem:fewer-arcs}
\cited{the fourth paper's Theorem~H~\cite[Theorem~6.5]{FourthPaper2026}, applied to the identity map; Lean-verified
there (\lean{MotionDepthCodes.RateBound.of_hom}).}
If $N(v)\subseteq N'(v)$ for every $v$ and $N'$ satisfies \eqref{eq:star}, then so does $N$.
\end{lemma}

Indeed, a polymatroid valid for $N$ is valid for $N'$, by diminishing returns. The last tool is contraction.

\begin{lemma}[Contraction]\label{lem:contraction}
\formal{Decomposition.lean}{contr, contr_isPolymatroid, contr_insert_eq}
Let $h$ be a polymatroid on $V$ and $U\subseteq V$. Then $h_U(X)=h(X\cup U)-h(U)$ is a polymatroid. If $h$ is valid
for $N$ and $N(v)\subseteq U\cup X$, then $h_U(X\cup\{v\})=h_U(X)$; in particular the restriction of $h_U$ to
$V\setminus U$ is valid for $N-U$. Moreover $h_U(w)=h(w)-I(w;U)$ for $w\notin U$.
\end{lemma}

\begin{outline}
The polymatroid axioms for $h_U$ follow from those for $h$, because $X\mapsto X\cup U$ preserves unions,
intersections and inclusions. If $N(v)\subseteq U\cup X$, the closure lemma (\cref{lem:closure}) gives
$h(X\cup U\cup\{v\})=h(X\cup U)$. The last identity is the definition of $I(w;U)$.
\end{outline}

\subsection{Formal status of the results}\label{sec:prelim-lean}

Every numbered result names its formal status in one of four ways. A line
\begin{center}
{\small\textsf{Lean:} \texttt{declaration, \dots}\hfill[\texttt{file}]}
\end{center}
gives the Lean declarations of a fully formalized result. The declarations live in the namespace \texttt{HalfRate}
and the files in the directory \leanfile{lean/HalfRate/}. The text after such a statement, labelled \emph{Outline of the
formal proof}, is a guide to the formal proof; it is a convenience, not a proof. A line beginning \textsf{Lean
(partial)} marks a result whose formal counterpart covers all but an explicitly named part, which the text proves.
The words \textsf{Not formalized} mark a result with no formal counterpart; it comes with a complete proof, labelled
\emph{Proof (not formalized)}. A line beginning \textsf{Cited} marks a result of an earlier paper of the series that is
used without proof, and says where it is proved and whether it is Lean-verified there. Where the English statement
simplifies, the Lean statement is authoritative. Numerical observations are labelled as such and are not used in
proofs. Numbered definitions carry a status line only when they are formalized.

The library depends on Lean~4 and Mathlib (both v4.34.1), on the Lean library of the fourth paper (whose
polymatroids, storage codes and rate bounds we use unchanged) and on that of its circular memo~\cite{CircularMemo2026},
imported through relative paths. It contains no \texttt{sorry}, no \texttt{native\_decide} and no unsafe evaluation,
and every declaration cited here depends only on the axioms \texttt{propext}, \texttt{Classical.choice} and
\texttt{Quot.sound} (\cref{app:verify}). \Cref{app:index} lists every declaration of this library cited in the paper.

\section{Nested elimination and the decomposition theorem}\label{sec:engine}

This section proves the engine of the paper. A \emph{piece} $H$ is eliminated one vertex at a time. When a vertex is
eliminated it must be recovered from the vertices eliminated before it, from one \emph{record} other than its own, and
possibly from one vertex outside the piece. Lemma~E bounds $2h(H)$ by the records plus the information that the outside
vertices have about $H$, provided the steps charged to the same outside vertex are nested. Contracting $H$ then reduces
the half bound to the rest of the relation, and Theorem~D composes pieces. Every other result of the paper is an
instance: it only has to exhibit pieces.

\subsection{Lemma E}\label{sec:lemmaE}

\begin{definition}[Elimination data]\label{def:elim-data}
\formal{NestedElimination.lean}{defect, totalDefect, Charged}
Let $H\subseteq V$. \emph{Elimination data} for $H$ consist of finitely many \emph{records} $P_i\subseteq H$, $i\in I$,
not necessarily disjoint; an \emph{elimination list} $v_0,\dots,v_{k-1}$, a listing of $H$ without repetitions (the
formal statement allows repetitions, which do not help); and maps $a,b\colon H\to I$ with $v\in P_{a(v)}$ and $b(v)\ne a(v)$ for every $v\in H$, the \emph{own} and the
\emph{supplying} record of $v$. At step $t$ the \emph{revealed set} is $S_t=\{v_0,\dots,v_{t-1}\}$ and the
\emph{supply} is $Q_t=S_t\cup P_{b(v_t)}$. For a polymatroid $g$ on $V$, the \emph{defect} of step $t$ is
$g(v_t\mid Q_t)$.

Let $W\subseteq V$. A \emph{charging} into $W$ assigns to every vertex of $H$, and so to its step, either nothing or a
vertex $w\in W$, such that the defect of a step with nothing vanishes (the step is \emph{free}) and a step assigned to $w$ satisfies
$g(v_t\mid Q_t\cup\{w\})=0$ (it is \emph{charged to $w$}). The charging is \emph{nested} if, whenever $t'<t$ are
consecutive steps charged to the same vertex, $Q_{t'}\cup\{v_{t'}\}\subseteq Q_t$.
\end{definition}

Consecutive nesting implies $Q_{t'}\cup\{v_{t'}\}\subseteq Q_t$ for all steps $t'<t$ charged to the same vertex,
since supplies grow along the chain.

\begin{lemma}[Lemma E: nested elimination]\label{lem:E}
\formal{NestedElimination.lean}{record_increment, elim_defect, mi_mono, defect_eq_mi, totalDefect_le, nested_elimination}
Let $g$ be a polymatroid on $V$ and fix elimination data for $H\subseteq V$.
\begin{enumerate}[label=(\alph*)]
\item\label{it:E0} (Defect form.) $2g(H)\le\sum_ig(P_i)+\sum_tg(v_t\mid Q_t)$.
\item\label{it:E1} (Charging.) If a charging into $W\subseteq V$ is nested, then
$\sum_tg(v_t\mid Q_t)\le\sum_{w\in W}I(w;H)$.
\item\label{it:E} (Lemma E.) Consequently, if a charging into $W$ is nested, then
\[
  2g(H)\;\le\;\sum_ig(P_i)+\sum_{w\in W}I(w;H).
\]
\end{enumerate}
\end{lemma}

\begin{outline}
\ref{it:E0} Let $m=\card I$; it is at least $2$ as soon as $H\ne\emptyset$, since $a(v)\ne b(v)$. Induction along the
list proves, for every $t$,
\[
  2g(H)+(m-2)\,g(S_t)\;\le\;\sum_ig(S_t\cup P_i)+\sum_{s\ge t}g(v_s\mid Q_s).
\]
At the end, $S_k=H$ contains every record and both sides equal $m\,g(H)$. In the inductive step $v=v_t$ is revealed: the
own record does not change, because $v\in P_{a(v)}$; the supplying record grows by
exactly the defect $g(v\mid S_t\cup P_{b(v)})$; and each of the other $m-2$ records grows by at most
$g(S_t\cup\{v\})-g(S_t)$, by submodularity (\lean{record_increment}). At $t=0$ this is~\ref{it:E0}.

\ref{it:E1} For $w\in W$ let $L(w)$ be the supply of the last step charged to $w$ together with its vertex, initially
$\emptyset$. Induction along the list proves
$\sum_{s\ge t}g(v_s\mid Q_s)\le\sum_{w\in W}\bigl(I(w;H)-I(w;L(w))\bigr)$, where $L$ is taken at step $t$. A free step
contributes nothing. For a step charged to $w$, the chain rule and $g(v_t\mid Q_t\cup\{w\})=0$ give
$g(v_t\mid Q_t)=I(w;Q_t\cup\{v_t\})-I(w;Q_t)$ (\lean{defect_eq_mi}); nesting gives $L(w)\subseteq Q_t$, and
$I(w;\cdot)$ is monotone (\lean{mi_mono}), so the defect is at most $I(w;Q_t\cup\{v_t\})-I(w;L(w))$, after which
$L(w)$ becomes $Q_t\cup\{v_t\}\subseteq H$. At the start $I(w;\emptyset)=0$.

\ref{it:E} Combine~\ref{it:E0} and~\ref{it:E1}.
\end{outline}

With $W=\emptyset$ and all steps free, part~\ref{it:E} is the inequality form of the exact partition identity of the
circular memo~\cite[Theorem~2.1]{CircularMemo2026}, which assumes that the records partition the ground set and
identifies the slack exactly; part~\ref{it:E0} is the inequality form of that memo's remark that dropping the
supplying condition subtracts the defects. What is new is part~\ref{it:E1}: a step may use one vertex outside the
piece, and several steps may use the same outside vertex, at the total price $I(w;H)$, provided their supplies are
nested. Validity of $g$ is not assumed, and the records need not be disjoint; validity enters in Theorem~D through the
closure lemma, and disjointness through subadditivity.

\paragraph{Why it matters.} Lemma~E is the only entropy inequality of the paper; everything after it is combinatorics.
For~(T) this means that the half-rate bounds needed by the impossibility programme of the fourth
paper~\cite{FourthPaper2026} can be proved by exhibiting finite combinatorial certificates. The charging of
part~\ref{it:E1} is what directed side information requires: in the directed classes of \cref{sec:directed}, an outside
vertex can be a parent of several vertices of a piece. Not every natural piece admits such a charging; \cref{rem:X15}
exhibits a minimum nonempty core that admits no piece certificate, with any number of records.

\subsection{Piece certificates and Theorem D}\label{sec:thmD}

\begin{definition}[Piece certificates and decompositions]\label{def:piece}
\formal{Decomposition.lean}{RelCharged, PieceCert, Decomp}
Let $N$ be a recovery relation on $V$, let $U\subseteq V$ be a \emph{context} and $H\subseteq V\setminus U$ a
\emph{piece}, and call the vertices of $V\setminus(U\cup H)$ \emph{outside}. A \emph{piece certificate} for $H$ in
context $U$ consists of elimination data for $H$ whose records $P_1,\dots,P_m$ are pairwise disjoint with union $H$,
and of an optional \emph{outside parent} $p(v)$, an outside vertex, for each $v\in H$, such that for every step $t$:
\begin{enumerate}[label=(C\arabic*),leftmargin=3em]
\item\label{it:C1} $N(v_t)\subseteq U\cup Q_t\cup\{p(v_t)\}$, where $\{p(v_t)\}=\emptyset$ if $v_t$ has no outside
parent;
\item\label{it:C2} if $t'<t$ are consecutive steps with the same outside parent, then
$Q_{t'}\cup\{v_{t'}\}\subseteq Q_t$.
\end{enumerate}
So every parent of $v_t$ outside the context has been eliminated, lies in the supplying record, or is the single
outside parent of $v_t$. A \emph{decomposition of $N$ from context $U$} is a list of pieces $H_1,\dots,H_k$ such that,
with $U_j=U\cup H_1\cup\dots\cup H_{j-1}$, each $H_j$ is disjoint from $U_j$ and has a piece certificate in context
$U_j$, and $U_{k+1}=V$. A \emph{decomposition} of $N$ is a decomposition from context $\emptyset$.
\end{definition}

\begin{theorem}[Theorem D: decomposition]\label{thm:D}
\formal{Decomposition.lean}{contr_isPolymatroid, RelCharged.charged, PieceCert.ineq, decomp_ineq, rateBound_of_decomp, card_sq_le_of_decomp}
Let $h$ be a polymatroid valid for $N$.
\begin{enumerate}[label=(\alph*)]
\item\label{it:D-piece} (Pieces.) If $H$ has a piece certificate in context $U$, then, with the contraction $h_U$ of
\cref{lem:contraction} and the information $I_{h_U}$ computed in $h_U$,
\[
  2h_U(H)\;\le\;\sum_{v\in H}h_U(v)+\sum_{w\notin U\cup H}I_{h_U}(w;H).
\]
\item\label{it:D-rel} (Relative form.) If $N$ has a decomposition from context $U$, then
\[
  2\bigl(h(V)-h(U)\bigr)\;\le\;\sum_{v\notin U}\bigl(h(U\cup\{v\})-h(U)\bigr).
\]
\item\label{it:D-main} (Theorem D.) If $N$ has a decomposition, then $N$ satisfies \eqref{eq:star}, and every storage
code $C\subseteq\alpha^V$ for $N$ over a nonempty finite alphabet satisfies $\card C^2\le\card\alpha^{\card V}$.
\end{enumerate}
\end{theorem}

\begin{outline}
\ref{it:D-piece} Apply Lemma~E to $g=h_U$, a polymatroid by \cref{lem:contraction}, with $W$ the set of outside
vertices. Condition~\ref{it:C1} and \cref{lem:contraction} turn each step into a free step or a step charged to its
outside parent (\lean{RelCharged.charged}), and \ref{it:C2} is the nesting. Finally
$\sum_ih_U(P_i)\le\sum_{v\in H}h_U(v)$ by subadditivity over the disjoint records.

\ref{it:D-rel} Induction over the list of pieces. For the first piece $H$ in context $U$, part~\ref{it:D-piece} applies;
the induction hypothesis, for the remaining pieces from context $U\cup H$, gives
$2(h(V)-h(U\cup H))\le\sum_{w\notin U\cup H}(h(U\cup H\cup\{w\})-h(U\cup H))$. Since
$h(U\cup H\cup\{w\})-h(U\cup H)=h_U(w)-I_{h_U}(w;H)$, the information terms cancel when the two inequalities are added,
and $h(V)-h(U)=h_U(H)+(h(V)-h(U\cup H))$.

\ref{it:D-main} Take $U=\emptyset$ in~\ref{it:D-rel}; the code bound is \cref{thm:codes}.
\end{outline}

The pieces are certified in the relation that remains after contracting the earlier pieces: parents in earlier
pieces are free. The bookkeeping is exact, in the sense that the piece inequalities add up to
$\sum_vh(v)-2h(V)$ with all outside charges cancelling; unfolding a decomposition therefore gives an explicit Shannon
certificate for \eqref{eq:star}, a nonnegative combination of submodularity and validity brackets.

\paragraph{Why it matters.} Theorem~D turns each class theorem below into a statement that decompositions exist, and
it is the only route by which this paper proves half-rate bounds. For~(T), a proof of the fourth paper's crux along
these lines would have to decompose the cone relations of oblivious schedules into certified pieces;
\cref{sec:limits} shows where the natural rules for choosing pieces break down.

\subsection{Building certificates}\label{sec:building}

Three tools produce piece certificates. The first removes the need to check nesting.

\begin{lemma}[Co-supply implies nesting]\label{lem:cosupply}
\formal{Positional.lean}{rev, relCharged_of_cosupply}
Let elimination data for $H$, with records partitioning $H$ and optional outside parents, satisfy~\ref{it:C1} at every
step. If all steps with the same outside parent have the same supplying record (the outside parents are
\emph{co-supplied}), then~\ref{it:C2} holds. In particular~\ref{it:C2} holds if every outside vertex is the outside
parent of at most one vertex of $H$ (the outside parents are \emph{private}).
\end{lemma}

\begin{outline}
For steps $t'<t$ with the common supplying record $P_b$,
$Q_{t'}\cup\{v_{t'}\}=S_{t'+1}\cup P_b\subseteq S_t\cup P_b=Q_t$. The formal statement takes the step conditions
position by position along the list, with \lean{rev} the revealed set before a position.
\end{outline}

\begin{lemma}[Fewer arcs and greedy certificates]\label{lem:greedy}
\formal{Pieces.lean}{RelCharged.mono, PieceCert.mono, exists_relCharged, pieceCert_of_elim}
Let $U\subseteq V$ and $H\subseteq V\setminus U$.
\begin{enumerate}[label=(\alph*)]
\item\label{it:mono} If $N(v)\subseteq N'(v)$ for every $v\in H$, every piece certificate for $H$ in context $U$ for
$N'$ is one for $N$.
\item\label{it:greedy} Let records $P_1,\dots,P_m$ partition $H$, with own and supplying records $a(v)\ne b(v)$ and
optional co-supplied outside parents $p(v)$. If every proper subset $S\subsetneq H$ misses a vertex $v\in H$ with
$N(v)\subseteq U\cup S\cup P_{b(v)}\cup\{p(v)\}$, then $H$ has a piece certificate in context $U$.
\end{enumerate}
\end{lemma}

\begin{outline}
\ref{it:mono} Condition~\ref{it:C1} only becomes weaker, and~\ref{it:C2} does not mention $N$.
\ref{it:greedy} Starting from $S=\emptyset$, repeatedly eliminate a vertex given by the hypothesis. Each step
satisfies~\ref{it:C1} by construction, and co-supply gives~\ref{it:C2} (\cref{lem:cosupply}). The list ends when $S=H$.
\end{outline}

\begin{lemma}[Lemma $\Bdir$: relative $2$-dicolourings]\label{lem:Bdir}
\formal{Pieces.lean}{bdirPiece}
Let $U\subseteq V$ and let $H=A\sqcup B\subseteq V\setminus U$, where $A$ and $B$ are acyclic for $N$. Suppose that
every vertex of $H$ has at most one parent outside $U\cup H$, and that every outside vertex has all its children in $H$
inside $A$, or all inside $B$. Then $H$ has a piece certificate in context $U$ (the formal proof uses the two records
$A$ and $B$).
\end{lemma}

\begin{outline}
The supplying record of a vertex of $A$ is $B$, and that of a vertex of $B$ is $A$; the outside parent of a vertex is
its parent outside $U\cup H$, if any. Apply \cref{lem:greedy}\ref{it:greedy}. Given a proper subset $S$ of $H$: if
$A\not\subseteq S$, acyclicity gives a vertex of $A\setminus S$ with no parent in $A\setminus S$, whose parents in $H$
are therefore revealed or in $B$; otherwise $A\subseteq S$, and a vertex of $B\setminus S$ with no parent in
$B\setminus S$ has all its parents in $H$ revealed. Co-supply holds because the children of an outside vertex lie in
one part, which has one supplying record.
\end{outline}

\paragraph{Why it matters.} Relative $2$-dicolourings are the workhorse of \cref{sec:directed}: directed triangles,
shortest directed cycles of in-degree-two vertices and the final step of Theorem~T are pieces of this kind. They need no
degree bound and allow an outside vertex to have many children in a piece, as long as these lie on one side.

\subsection{Global certificates and pull-back}\label{sec:global}

A piece certificate with $H=V$ has no outside vertices. Its order-free form is the following.

\begin{definition}[Global elimination certificates]\label{def:elimcert}
\formal{Corollaries.lean}{ElimCert}
A \emph{global elimination certificate} for $N$ consists of records $P_1,\dots,P_m$ partitioning $V$ and own and
supplying records $a(v)\ne b(v)$ with $v\in P_{a(v)}$, such that every proper subset $S\subsetneq V$ misses a vertex
$v$ with $N(v)\subseteq S\cup P_{b(v)}$.
\end{definition}

Equivalently, there is an elimination list in which the parents of each vertex that are not earlier in the list lie in
its supplying record: given the list and a proper subset $S$, the first vertex of the list outside $S$ is eligible.

\begin{proposition}[Global certificates]\label{prop:global}
\formal{Corollaries.lean}{ElimCert.rateBound}
\formal{NestedElimination.lean}{sum_records_step}
If $N$ has a global elimination certificate, then $N$ satisfies \eqref{eq:star}.
\end{proposition}

\begin{outline}
Downward induction on $\card{V\setminus S}$ proves $2h(V)+(m-2)h(S)\le\sum_ih(S\cup P_i)$, as in
\cref{lem:E}\ref{it:E0}, with every defect zero by validity and the closure lemma; the inductive step is
\lean{sum_records_step}. At $S=\emptyset$, subadditivity over the records gives \eqref{eq:star}.
\end{outline}

With three records and cyclic supply, $b(v)$ the record after $a(v)$, this is the three-record criterion of the
circular memo~\cite[Lemma~2.2]{CircularMemo2026} for partitions; the memo also allows overlapping records. Global
certificates pull back along homomorphisms of recovery relations.

\begin{lemma}[Pull-back]\label{lem:pullback}
\formal{Corollaries.lean}{ElimCert.pullback}
Let $N$ and $N'$ be recovery relations on $V$ and $W$, and let $\varphi\colon V\to W$ satisfy $\varphi(u)\in
N'(\varphi(v))$ whenever $u\in N(v)$. If $N'$ has a global elimination certificate with records $P'_i$ and maps
$a',b'$, then $N$ has one with records $\varphi^{-1}(P'_i)$ and maps $a'\circ\varphi$, $b'\circ\varphi$.
\end{lemma}

\begin{outline}
Given a proper subset $S\subsetneq V$, let $S'$ be the set of $x\in W$ whose whole fibre $\varphi^{-1}(x)$ lies in $S$.
It is a proper subset of $W$, so some $x\notin S'$ satisfies $N'(x)\subseteq S'\cup P'_{b'(x)}$; choose
$v\in\varphi^{-1}(x)\setminus S$. A parent $u$ of $v$ has $\varphi(u)\in N'(x)$: either $\varphi(u)\in S'$, so $u\in S$,
or $\varphi(u)\in P'_{b'(x)}$, so $u$ lies in the supplying record of $v$.
\end{outline}

No looplessness is needed: a relation with a global certificate is loopless, since a loop at $v$ would put $v$ in its
own supplying record. The fourth paper's Theorem~H~\cite[Theorem~6.5]{FourthPaper2026} pushes polymatroids forward
along homomorphisms and so transfers \eqref{eq:star} itself; \cref{lem:pullback} transfers the certificate instead,
and needs no polymatroid argument.

\paragraph{Why it matters.} One certificate on a small target proves \eqref{eq:star} for every relation that maps to
it, whatever its size or degree. For~(T), an impossibility proof could map the cone relations of a schedule to certified
targets; a construction must avoid side information that admits such a map. Global certificates are weaker than
decompositions (\cref{rem:petersen}), which is why the classes of \cref{sec:undirected,sec:directed} use Theorem~D.

\section{Base cases and symmetric side information}\label{sec:undirected}

This section treats the base case of $2$-dicolourable relations and two classes of graphs: graphs of circular chromatic
number below three, at any degree, and triangle-free graphs of maximum degree three.

\begin{proposition}[$2$-dicolourable relations]\label{prop:twoAcyclic}
\formal{Corollaries.lean}{IsAcyclicIn, rateBound_of_twoAcyclic}
If $V$ is the disjoint union of two sets $A$ and $B$ that are acyclic for $N$, then $N$ satisfies \eqref{eq:star}.
\end{proposition}

\begin{outline}
A global elimination certificate with records $A$ and $B$, each the supplying record of the other: while $A$ is not
revealed, eliminate a vertex of the unrevealed part of $A$ with no parent there; then do the same in $B$. Apply
\cref{prop:global}.
\end{outline}

Bipartite graphs are examples, and so is every loopless relation whose underlying graph is a forest: its directed cycles are
digons, so a proper $2$-colouring of the underlying graph is a $2$-dicolouring. The bound is classical: for a valid $h$,
each of $A$ and $B$ determines the other in topological order, so $h(V)=h(A)=h(B)$.

\paragraph{Why it matters.} This is the base case of every decomposition below: pieces end, after their cyclic parts have
been removed, in $2$-dicolourings. For~(T) it says that side information without odd dependency structure, in this
directed sense, gives no coding gain at all.

\subsection{Circular chromatic number below three}\label{sec:circular}

The \emph{circular clique} $K_{p/q}$ has vertices $\Z/p$, with $i\sim j$ if and only if $q\le(j-i)\bmod p\le p-q$.
The \emph{Andrásfai graph} $\Andr(k)$, $k\ge2$, is $K_{(3k-1)/k}$. The circular chromatic number $\chi_c(G)$ of a finite
graph is the least $p/q$ such that $G$ maps homomorphically to $K_{p/q}$~\cite{Vince1988,Zhu2001}; so $\chi_c(G)<3$ if
and only if $G$ maps to some $K_{p/q}$ with $p<3q$.

\begin{theorem}[Circular chromatic number below three]\label{thm:circular}
\formal{Corollaries.lean}{andrasfaiCert, rateBound_of_circularHom'}
Let $G$ be a finite simple graph with a homomorphism to a circular clique $K_{p/q}$ with $q>0$ and $p<3q$, that is,
with $\chi_c(G)<3$. Then $N_G$ satisfies \eqref{eq:star}.
\end{theorem}

This is the circular memo's theorem~\cite[Corollary~3.2]{CircularMemo2026}, Lean-verified there through the fourth
paper's Theorem~H. We re-derive it as a direct instance of global certificates.

\begin{outline}
The Andrásfai graph $\Andr(k)$ has a global elimination certificate with three records
(\lean{andrasfaiCert}): the intervals $[0,k)$, $[k,2k)$ and $[2k,3k-1)$, each supplied by the next one cyclically.
The elimination order is the circular memo's interleaved order, with keys $3i$, $3j+1$ and $3\ell+2$ for the $i$-th,
$j$-th and $\ell$-th vertex of the three intervals, and the memo's adjacency lemma (\lean{CircularMemo.neighbor_order})
says that every later neighbour of a vertex lies in the next interval. For $k=q+1$ we have $pk\le(3k-1)q$, so the
Bondy--Hell map $i\mapsto\lfloor(3k-1)i/p\rfloor$ is a homomorphism $K_{p/q}\to\Andr(k)$~\cite{BondyHell1990}, formalized
in the fourth paper (\lean{MotionDepthCodes.circularCliqueHom}). Compose it with the given homomorphism, pull the
certificate back (\cref{lem:pullback}) and apply \cref{prop:global}.
\end{outline}

\paragraph{Why it matters.} The theorem has no degree assumption: every graph of circular chromatic number below three,
for instance every odd cycle of length at least five and every graph mapping to one, satisfies the half bound for all
codes over all alphabets. For~(T), symmetric side information of this kind cannot carry a coding gain however dense it
is. The bound does not reach $\chi_c=3$; the Petersen graph, which has $\chi_c=3$~\cite{JHS1997}, is covered by the
fourth paper's certificate~\cite[Corollary~6.9]{FourthPaper2026} and by the next theorem.

\subsection{Triangle-free graphs of maximum degree three}\label{sec:N3}

\begin{theorem}[Theorem N3]\label{thm:N3}
\formal{N3.lean}{CycIn, CycIn.chordless, CycIn.share, cyclePiece, decomp_of_cycles, n3_decomp, rateBound_n3, card_sq_le_n3}
Let $G$ be a triangle-free graph of maximum degree at most~$3$. Then $N_G$ has a decomposition. Hence $N_G$ satisfies
\eqref{eq:star}, and every storage code $C$ for $G$ over a nonempty finite alphabet satisfies
$\card C^2\le\card\alpha^{\card V}$.
\end{theorem}

\begin{outline}
By strong induction on the set $R$ of remaining vertices, with context $V\setminus R$, \lean{decomp_of_cycles}
constructs a decomposition from every context. If $G[R]$ contains a cycle, let $C=c_0c_1\cdots c_{g-1}$ be a shortest one
(\lean{CycIn}); $g\ge4$ because $G$ is triangle-free. Three facts:
\begin{enumerate}[label=(\roman*)]
\item $C$ is chordless (\lean{CycIn.chordless}), since a chord would close a shorter cycle;
\item if a vertex $w\in R$ off $C$ is adjacent to $c_i$ and $c_j$, $i<j$, it closes cycles of lengths $j-i+2$ and
$g-(j-i)+2$, both at least $g$; hence $g=4$ and $c_i,c_j$ are opposite (\lean{CycIn.share});
\item every $c_i$ has at most one neighbour off $C$, since it has two on $C$ and degree at most three.
\end{enumerate}
The piece certificate of $C$ in context $V\setminus R$ (\lean{cyclePiece}) eliminates $c_0,\dots,c_{g-1}$ in order with
three records given by colours: $c_0$ has colour~$0$ and $c_{g-1}$ colour~$1$; if $g$ is odd, $c_{g-2}$ has
colour~$2$; every other $c_i$ has colour~$1$ for odd $i$ and~$0$ for even $i$. The supplying record of $c_t$ is the
colour of $c_{t+1}$, and that of $c_{g-1}$ is colour~$0$. At a step $t\ge1$ the cycle neighbours of $c_t$ are
$c_{t-1}$, already eliminated, and $c_{t+1}$, in the supplying record; at step~$0$ they are $c_1$ and $c_{g-1}$, both of
colour~$1$. The remaining neighbour of $c_t$, if it lies in $R\setminus C$, is its outside parent. By~(ii) outside
parents are private, except for opposite vertices of a $4$-cycle, which have equal supplying records ($c_0,c_2$ are
supplied by colour~$1$, and $c_1,c_3$ by colour~$0$); \cref{lem:cosupply} applies. If $G[R]$ is a forest, a proper
$2$-colouring of it is a piece by \cref{lem:Bdir}. Theorem~D concludes.
\end{outline}

By \cref{lem:fewer-arcs}, the half bound holds also for every relation $N$ with $N(v)\subseteq N_G(v)$ for all $v$,
that is, for every relation whose arcs join adjacent vertices of $G$, in one direction or in both;
\lean{decomp_of_cycles} gives decompositions for such relations as well, its hypothesis on $3$-cycles being vacuous here.
Several known results are special cases.
\begin{itemize}[leftmargin=2em]
\item The fourth paper's Theorem~S~\cite[Theorem~6.13]{FourthPaper2026} proves the half bound for scalar linear codes,
over every field, on a directed class that contains these graphs; Theorem~N3 is its all-code version on triangle-free
subcubic graphs. For scalar linear codes the exact value is $\card V-\nu(G)$ for the minrank, with the integral matching
number $\nu$~\cite[Theorem~6.16]{FourthPaper2026}; for all codes the guessing number is the fractional matching number
$\nustar(G)$ (\cref{thm:gn}). The two differ already on $C_5$, where $\nu=2$ and $\nustar=5/2$.
\item The subcubic memo~\cite[Theorem~7.1]{SubcubicMemo2026} proves the half bound, in Lean, for triangle-free subcubic
graphs all of whose edges lie in four-cycles, through a homomorphism to $K_{8/3}$; this is a special case. Its
information-set theorem~\cite[Theorem~2.3]{SubcubicMemo2026} bounds codes that are reversible on the rows by
$q^{\nu(G)}$, which is stronger on that class of codes.
\item The circular memo's theorem (\cref{thm:circular}) covers graphs of every degree with $\chi_c<3$, and the fourth
paper's certificates cover the graphs that map to the Petersen graph~\cite[Corollary~6.9]{FourthPaper2026}. Theorem~N3
needs no target: it covers every triangle-free subcubic graph. In particular it proves unconditionally the fourth paper's
conditional result for cubic graphs of large girth~\cite[Corollary~6.10]{FourthPaper2026}, which assumed that such graphs
map to $C_5$ or to $K_{14/5}$.
\item Dang~\cite{Dang2015} computes the guessing numbers of several families of planar triangle-free graphs and
conjectures that the fractional clique cover strategy is optimal on every planar triangle-free
graph~\cite[Conjecture~15]{Dang2015}. With \cref{thm:gn} below, Theorem~N3 proves this for maximum degree three,
without planarity.
\end{itemize}

\paragraph{Why it matters.} No storage code on a triangle-free subcubic graph beats the pairing, whatever its alphabet
and whether or not it is linear. With \cref{cor:paid}, this answers the fourth paper's undirected
crux~\cite[Question~6.1]{FourthPaper2026} affirmatively at maximum degree three: in a locally bipartite family of maximum
degree at most three, only $o(\card V)$ vertices lie on triangles, deleting them leaves a triangle-free subcubic graph, and
so every storage code has rate at most $\frac12+o(1)$. For~(T), symmetric subcubic side
information that is triangle-free cannot carry a coding gain. Native cone relations are directed, however, and their
symmetrizations can contain linearly many triangles~\cite[Proposition~3.10]{DirectedRecovery2026}; this is why
\cref{sec:directed} allows directed triangles.

\section{Directed relations without transitive triangles}\label{sec:directed}

By \cref{lem:tt-paired}, a loopless relation is TT-free exactly when every coordinate has paired distance at least
five. This section proves the half bound for three classes of such relations: subcubic underlying graphs
(Theorem~U3), relations whose vertices of in-degree at least three are acyclic (Theorem~T), and block structures with
private outside parents (Theorem~PC), which contain all TT-free Cayley relations $\Cay(G,\{s,s^{-1},t\})$
(Corollary~CC). Directed triangles are allowed throughout; they are the only triangles a TT-free relation has.

\begin{lemma}[Triangles of TT-free relations]\label{lem:dtri}
\formal{U3.lean}{dtri_orient}
Let $N$ be TT-free, and let $x_0,x_1,x_2$ be distinct and pairwise adjacent in the underlying graph, with
$x_0\to x_1$. Then $x_1\to x_2$ and $x_2\to x_0$, and none of the arcs $x_1\to x_0$, $x_2\to x_1$, $x_0\to x_2$ is
present.
\end{lemma}

\begin{outline}
If $x_2\to x_1$, then either $x_0\to x_2$, and $x_0\to x_2\to x_1$ with $x_0\to x_1$ is a transitive triangle, or
$x_2\to x_0$, and $x_2\to x_0\to x_1$ with $x_2\to x_1$ is one. So $x_1\to x_2$, and then $x_0\to x_2$ would close
the transitive triangle $x_0\to x_1\to x_2$; hence $x_2\to x_0$. The reverse arcs $x_1\to x_0$ and $x_2\to x_1$ would
close the transitive triangles $x_1\to x_2\to x_0$, $x_1\to x_0$ and $x_2\to x_0\to x_1$, $x_2\to x_1$.
\end{outline}

\subsection{Subcubic underlying graphs: Theorem U3}\label{sec:U3}

\begin{theorem}[Theorem U3]\label{thm:U3}
\formal{U3.lean}{und, sub_nbhd_und, dtriPiece, triPiece, rateBound_u3, card_sq_le_u3}
Let $N$ be loopless and TT-free, and let its underlying graph have maximum degree at most~$3$. Then $N$ satisfies
\eqref{eq:star}, and every storage code $C$ for $N$ over a nonempty finite alphabet satisfies
$\card C^2\le\card\alpha^{\card V}$.
\end{theorem}

\begin{outline}
Let $G=U(N)$; since $N$ is loopless, $N(v)\subseteq N_G(v)$ (\lean{sub_nbhd_und}). Run the shortest-cycle decomposition
of Theorem~N3 on $G$ (\lean{decomp_of_cycles}). Shortest cycles of length at least four are pieces for $N_G$ and hence
for $N$, which has fewer arcs (\cref{lem:greedy}\ref{it:mono}); the acyclic remainder is a relative $2$-colouring, built
directly for $N$ by \cref{lem:Bdir}. The new case is a shortest
cycle of length three. By \cref{lem:dtri} it is a directed $3$-cycle $c_0\to c_1\to c_2\to c_0$ without reverse arcs
(\lean{triPiece}). Take $A=\{c_0,c_2\}$, which carries the single arc $c_2\to c_0$ and is acyclic, and $B=\{c_1\}$
(\lean{dtriPiece}). Every $c_i$ has two neighbours on the triangle, hence at most one parent outside it. Two distinct
triangle vertices cannot share an outside parent $w$: with the arc between them, $w$ would complete a transitive
triangle. So outside parents are private, and \cref{lem:Bdir} applies.
\end{outline}

When the underlying graph is triangle-free, Theorem~U3 follows from Theorem~N3 and \cref{lem:fewer-arcs}; its new
content is directed triangles. The fourth paper's Theorem~S~\cite[Theorem~6.13]{FourthPaper2026} proves the half bound
for scalar linear codes on a larger class: its condition ``no recovery triangle'' is TT-freeness, and its reciprocity
condition asks that a vertex with three parents have no child that is not a parent, which holds when the underlying
graph is subcubic. Theorem~U3 is thus the all-code version of Theorem~S on relations with subcubic underlying graph;
the all-code version of Theorem~S on its whole class is open.

\paragraph{Why it matters.} Directed triangles do not help a code. This matters for~(T) because native cone relations
can contain linearly many disjoint directed $3$-cycles~\cite[Proposition~3.10]{DirectedRecovery2026}, which rule out the
symmetric hypothesis of Theorem~N3 but not the hypotheses of this section.

\subsection{Acyclic high in-degree: Theorem T}\label{sec:T}

\begin{theorem}[Theorem T]\label{thm:T}
\formal{TheoremT.lean}{DCyc, exists_hasDCyc_of_sourceless, DCyc.chordless, DCyc.acyclic_of_missing, singlePiece, i5Piece, block_decomp, rateBound_thmT, card_sq_le_thmT}
Let $N$ be loopless and TT-free, and suppose that the set $\{v:\card{N(v)}\ge3\}$ of vertices with at least three
parents is acyclic for $N$. Then $N$ satisfies \eqref{eq:star}, and every storage code $C$ for $N$ over a nonempty
finite alphabet satisfies $\card C^2\le\card\alpha^{\card V}$. No bound on in- or out-degrees is assumed.
\end{theorem}

\begin{outline}
By strong induction on the set $R$ of remaining vertices (\lean{block_decomp}, with one block), in-degrees being counted
inside $R$:
\begin{enumerate}[label=(\alph*)]
\item If some $v\in R$ has at most one parent in $R$, then $\{v\}$ is a piece, its remaining parent being its outside
parent (\lean{singlePiece}, \cref{lem:Bdir} with $A=\{v\}$ and $B=\emptyset$).
\item Otherwise let $W_R$ be the set of vertices with exactly two parents in $R$. If $W_R$ is not acyclic, it contains a
nonempty set in which every vertex has a parent, hence a directed cycle (\lean{exists_hasDCyc_of_sourceless}). A
shortest directed cycle $C$ in $W_R$ (\lean{DCyc}) is induced (\lean{DCyc.chordless}): the only parent of each
$c\in C$ on $C$ is its predecessor. Hence each $c\in C$ has exactly one further parent in $R$, off $C$, its outside
parent $p(c)$. Fix $c_0\in C$ and let $A=\{c\in C:p(c)=p(c_0)\}$ and $B=C\setminus A$. If $c_0\to c_1$ on $C$, then
$p(c_1)\ne p(c_0)$, since otherwise $p(c_0)\to c_0\to c_1$ with $p(c_0)\to c_1$ is a transitive triangle. So $A$ and
$B$ are proper subsets of an induced cycle and hence acyclic (\lean{DCyc.acyclic_of_missing}), and every outside vertex
has its children on $C$ inside $A$ (if it is $p(c_0)$) or inside $B$. \Cref{lem:Bdir} applies (\lean{i5Piece}).
\item Otherwise $W_R$ is acyclic, and so is the set of vertices with at least three parents in $R$, a subset of the
acyclic set of the hypothesis. These two sets partition $R$, which is a piece by \cref{lem:Bdir}.
\end{enumerate}
Theorem~D concludes.
\end{outline}

Strongly connected components can be treated separately, since a parent in an earlier component can be contracted. The
formal statement is a block form.

\begin{corollary}[Corollary $\Tint$]\label{cor:Tint}
\partformal{TheoremT.lean}{TTFreeOn, BlockOrder, blocks_decomp, rateBound_blocks}{the existence of a topological order of
the strongly connected components, which is standard; it is proved below.}
Let $N$ be loopless. Suppose that no strongly connected component of $N$ contains a transitive triangle, and that the set
$\Tint$ of vertices with at least three parents inside their own strongly connected component is acyclic. Then $N$
satisfies \eqref{eq:star}. There is no restriction on parents outside a vertex's component.
\end{corollary}

The formal theorem (\lean{rateBound_blocks}) is the following block form: if $N$ is loopless and $V$ is the union of
pairwise disjoint blocks
$S_1,\dots,S_r$ such that every parent of a vertex of $S_j$ lies in $S_1\cup\dots\cup S_j$, no block contains a
transitive triangle, and each set $\{v\in S_j:\card{N(v)\cap S_j}\ge3\}$ is acyclic, then $N$ satisfies
\eqref{eq:star}.

\begin{outline}
Induction over the list of blocks (\lean{blocks_decomp}): each block $S_j$ is decomposed by the procedure of Theorem~T
run inside $S_j$ (\lean{block_decomp}), with the earlier blocks in the context. Parents in earlier blocks are then free,
so in-degrees in~(a)--(c) are counted inside the remaining part of $S_j$, and TT-freeness is needed only inside $S_j$
(\lean{TTFreeOn}).
\end{outline}

\begin{informalproof}
We derive the corollary from the block form. The \emph{condensation} of $N$ has the strongly connected components as vertices and an arc $S\to S'$, $S\ne S'$,
whenever some arc of $N$ goes from $S$ to $S'$. It has no directed cycle, since components on such a cycle would be
strongly connected to each other. Hence it has a topological order: repeatedly remove a component that receives no arc
from the remaining ones. Listing the components $S_1,\dots,S_r$ in this order, every parent of a vertex of $S_j$ lies in
$S_1\cup\dots\cup S_j$. The set $\{v\in S_j:\card{N(v)\cap S_j}\ge3\}$ is $\Tint\cap S_j$, which is acyclic because
acyclicity passes to subsets. The block form applies.
\end{informalproof}

When $\Tint=\emptyset$, that is, when every vertex has at most two parents inside its component, Corollary~$\Tint$ is
the TT-free case of the two-input theorem of \emph{Directed Recovery
Structures}~\cite[Theorem~5.1]{DirectedRecovery2026}: TT-free components contain no bidirected triangles, so that
theorem's correction term vanishes. The corollary strictly contains this case, because it allows vertices with three or
more internal parents as long as they induce an acyclic sub-relation. Its proof dismantles a bidirected odd cycle by
contracting one digon, a shortest directed cycle of in-degree-two vertices, and does not use the directed Brooks
theorem on which the proof in~\cite{DirectedRecovery2026} relies. Conversely, \cite[Theorem~5.1]{DirectedRecovery2026}
also covers components that contain transitive triangles, which we do not.

\paragraph{Why it matters.} Theorem~T and Corollary~$\Tint$ have no degree bound. For~(T) they locate where a coding
gain on TT-free side information must come from: from directed cycles through vertices that have at least three parents
in their own component. In a cone relation, after the middle vertices of transitive triangles have been deleted at the
price of \cref{cor:paid}, a construction needs such cycles on a large fraction of the coordinates
(\cref{sec:construction}).

\subsection{Blocks with private outside parents: Theorem PC}\label{sec:PC}

The pieces of Theorems~N3, U3 and~T are small. The next theorem takes whole blocks as pieces, in any order; a block
then needs only an internal certificate.

\begin{definition}[Elimination certificates of a block]\label{def:elimon}
\formal{Blocks.lean}{ElimOn}
Let $K\subseteq V$. An \emph{elimination certificate} of $N[K]$ consists of an elimination list of $K$ without
repetitions, records partitioning $K$, and own and supplying records $a(v)\ne b(v)$ with $v\in P_{a(v)}$, such that at
every step $t$ the parents of $v_t$ inside $K$ lie in $S_t\cup P_{b(v_t)}$.
\end{definition}

This is a global elimination certificate of the induced relation $N[K]$ (\cref{def:elimcert}), presented by its list.

\begin{lemma}[Internal certificates give pieces]\label{lem:elimon-piece}
\formal{Blocks.lean}{ElimOn.piece}
Let $N[K]$ have an elimination certificate. Suppose that every vertex of $K$ has at most one parent outside $K$, and that
every vertex outside $K$ is a parent of at most one vertex of $K$. Then $K$ has a piece certificate in every context $U$
disjoint from $K$.
\end{lemma}

\begin{outline}
Keep the list and the records. The outside parent of $v\in K$ is its parent outside $U\cup K$, if any; its other
parents lie in $U$ or in $K$, where the internal certificate applies. Outside parents are private, so
\cref{lem:cosupply} gives nesting. (The formal statement asks for at most one parent in $V\setminus(U\cup K)$ only.)
\end{outline}

\begin{lemma}[Positional certificates; chordless digon cycles]\label{lem:cycle-cert}
\formal{Circulant.lean}{elimOn_ofFn, elimOn_cycle}
Let $K=\{k_0,\dots,k_{m-1}\}$ be nonempty and listed without repetitions.
\begin{enumerate}[label=(\alph*)]
\item\label{it:positional} Let $\mathrm{col}$ and $\mathrm{bcol}$ be colourings of the positions $0,\dots,m-1$ with
$\mathrm{col}(j)\ne\mathrm{bcol}(j)$ for every $j$. If $k_i\to k_j$ implies $i<j$ or $\mathrm{col}(i)=\mathrm{bcol}(j)$,
then $N[K]$ has an elimination certificate (the formal proof uses the list $k_0,\dots,k_{m-1}$ and the colour classes as
records).
\item\label{it:cycle} If $m\ge4$ and every arc inside $K$ joins cyclically consecutive vertices ($k_i\to k_j$ only if
$i\equiv j\pm1\pmod m$), then $N[K]$ has an elimination certificate.
\end{enumerate}
\end{lemma}

\begin{outline}
\ref{it:positional} The own record of $k_j$ is its colour and its supplying record is $\mathrm{bcol}(j)$; a parent $k_i$
that comes later has colour $\mathrm{bcol}(j)$. \ref{it:cycle} Use the colouring of the cycle pieces of Theorem~N3, with
$\mathrm{bcol}(j)=\mathrm{col}(j+1)$ for $j<m-1$ and $\mathrm{bcol}(m-1)=0$.
\end{outline}

Along a list $k_0,\dots,k_{L-1}$, call an arc $k_u\to k_v$ with $u\ge v+3$ a \emph{downward chord}.

\begin{lemma}[Lemma S: chorded paths]\label{lem:S}
\formal{PathCert.lean}{DownChord, pathRoot, pathRoot_succ_ne, elimOn_path}
Let $K=\{k_0,\dots,k_{L-1}\}$ be nonempty and listed without repetitions. Suppose that every arc $k_i\to k_j$ inside
$K$ joins consecutive positions, $\card{i-j}=1$, or skips at least two positions, $\card{i-j}\ge3$, and that every vertex
of $K$ is the head of at most one downward chord and the tail of at most one. Then $N[K]$ has an elimination certificate
(the formal proof uses the list $k_0,\dots,k_{L-1}$ and, as records, the classes of the equivalence relation on positions
generated by the \emph{links} $u\sim v+1$, one for each downward chord $k_u\to k_v$, together with one empty
record).
\end{lemma}

\begin{outline}
When $k_j$ is eliminated, its parents inside $K$ that come later are $k_{j+1}$ (if $j<L-1$) and the tail $k_u$ of the
downward chord to $k_j$, if there is one; these two are linked, so they lie in one class, which serves as supplying record
(through \cref{lem:cycle-cert}\ref{it:positional}). It remains to show that $j$ and $j+1$ are never in the same class
(\lean{pathRoot_succ_ne}). A link joins $u$ to $v+1\le u-2$. A position is the upper end of at most one link, from the
downward chord out of it, and the lower end of at most one link, from the downward chord into the position below it,
because these chords are unique. So the classes are chains whose positions decrease by at least two at each link:
following links downwards from any position reaches a lowest position, its \emph{root} (\lean{pathRoot}), and two
positions of one class differ by at least two.
\end{outline}

A \emph{chorded digon path} is a set $K=\{k_0,\dots,k_{L-1}\}$ such that $k_i\leftrightarrow k_{i+1}$ is a digon for
$0\le i<L-1$ and the other arcs inside $K$, the \emph{chords}, number at most one into and at most one out of every vertex.
If $N$ is loopless and $N[K]$ is TT-free, a chorded digon path satisfies the hypotheses of Lemma~S
(\lean{ChordedPath.elimOn}): a chord between $k_i$ and $k_{i\pm2}$ would close a transitive triangle with the two digons
around $k_{i\pm1}$.

\begin{proposition}[Proposition CC: chorded circulant cycles]\label{prop:CC}
\formal{Circulant.lean}{circCol, circB, elimOn_circulant, circulant_sweep}
Let $3\le k\le m-3$ and let $K=\{k_0,\dots,k_{m-1}\}$ be listed without repetitions so that every arc inside $K$ has the
form $k_{j\pm1}\to k_j$ or $k_{j-k}\to k_j$, indices modulo $m$. Put $\delta=m-1-k$.
\begin{enumerate}[label=(\alph*)]
\item\label{it:sweep} If $\delta$ divides neither $m-1$ nor $m-3$, then $N[K]$ has an elimination certificate (the
formal proof uses the sweep along $k_0,k_1,\dots,k_{m-1}$ described below).
\item\label{it:arith} If $2k\not\equiv\pm1\pmod m$, then the condition of~\ref{it:sweep} holds for the offset $k$ or for
the offset $m-k$.
\end{enumerate}
\end{proposition}

\begin{outline}
\ref{it:sweep} (\lean{elimOn_circulant}, through \cref{lem:cycle-cert}\ref{it:positional}.) Position $0$ has a record of
its own; a position $p\ne0$ gets the record of its residue modulo $\delta$, except that the residue class of $m-1$ is merged
into that of $1$ (\lean{circCol}). The supplying record of $p<m-1$ is the record of $p+1$, and that of $m-1$ is the record
of position~$0$ (\lean{circB}). The later parents of $k_p$ are $k_{p+1}$; for $p=0$ also $k_{m-1}$, which has the merged
colour of $k_1$; and, for $p<k$, the tail $k_{p-k+m}$ of the chord, whose position $p-k+m=(p+1)+\delta$ is congruent to
$p+1$. Own and supplying records differ: for $1\le p<m-1$ the residues of $p$ and $p+1$ differ, and the merge identifies
them only if $p\equiv m-1$ and $p+1\equiv1$, which forces $\delta\mid m-1$, or $p\equiv1$ and $p+1\equiv m-1$, which
forces $\delta\mid m-3$.
\ref{it:arith} (\lean{circulant_sweep}.) For the offset $m-k$ the role of $\delta$ is played by $\delta'=k-1$, and
$m-1=\delta+\delta'+1$, $m-3=\delta+\delta'-1$. If the condition fails for both offsets, then $\delta\mid\delta'\pm1$ and
$\delta'\mid\delta\pm1$ with $\delta,\delta'\ge2$, which forces $\card{\delta-\delta'}=1$ and then $2k=m\pm1$.
\end{outline}

\emph{With the mirror sweep, if $2k\not\equiv\pm1\pmod m$, then $N[K]$ has an elimination certificate.} Indeed, by
part~\ref{it:arith} the condition of part~\ref{it:sweep} holds for the offset $k$, and part~\ref{it:sweep} applies, or
for the offset $m-k$. In the second case relabel $k'_j=k_{-j}$: the relabelling $j\mapsto-j$ turns the arcs
$k_{j\pm1}\to k_j$ into arcs of the same form and the chords $k_{j-k}\to k_j$ into chords $k'_{j-(m-k)}\to k'_j$ of
offset $m-k$, so part~\ref{it:sweep} applies to the list $k'_0,k'_1,\dots,k'_{m-1}$, that is, to $k_0,k_{m-1},\dots,k_1$.
This is the \emph{mirror sweep}; it is not a single Lean declaration, and the formal proof of Corollary~CC below realizes
it by indexing cosets with powers of $s^{-1}$. The condition $2k\not\equiv0$, which makes the chords one-way, is not
needed. In a TT-free circulant $\Cay(\Z_m,\{1,-1,k\})$ with one-way chords the hypotheses
hold, since chords of length $2$ and chords with $2k\equiv\pm1$ close transitive triangles.

\begin{theorem}[Theorem PC: blocks with private outside parents]\label{thm:PC}
\formal{PathCert.lean}{ChordedPath, ChordedPath.elimOn, rateBound_pc}
\formal{Blocks.lean}{ChordlessSubcubic, localGraph, chordless_block_decomp}
\formal{N3.lean}{block_decomp_of_cycles}
Let $N$ be loopless and let $V$ be partitioned into blocks. Suppose:
\begin{enumerate}[label=(P\arabic*),leftmargin=3em]
\item\label{it:P1} every vertex has at most one parent outside its block;
\item\label{it:P2} every vertex is a parent of at most one vertex of each other block;
\item\label{it:P3} every block $K$ is
  \begin{enumerate}[label=(\alph*)]
  \item \emph{chordless subcubic}: every arc inside $K$ belongs to a digon, every vertex of $K$ has at most three
  parents, and $N[K]$ is TT-free; or
  \item a chorded digon path with $N[K]$ TT-free; or
  \item such that $N[K]$ has an elimination certificate.
  \end{enumerate}
\end{enumerate}
Then $N$ satisfies \eqref{eq:star}; consequently every storage code $C$ for $N$ over a nonempty finite alphabet satisfies
$\card C^2\le\card\alpha^{\card V}$ (\cref{thm:codes}). Out-degrees are not restricted.
\end{theorem}

\begin{outline}
By Lemma~S, blocks of kind~(b) have elimination certificates and are therefore of kind~(c)
(\lean{ChordedPath.elimOn}). The blocks are taken as pieces in an arbitrary order (\lean{rateBound_pc}), each with the earlier blocks as context.
By~\ref{it:P1} and~\ref{it:P2}, every vertex of a block $K$ has at most one parent outside $K$, and outside parents are
private. A block of kind~(c) is a piece by \cref{lem:elimon-piece}. A block of kind~(a) is decomposed relative to its
context by the shortest-cycle procedure of Theorem~N3 (\lean{chordless_block_decomp}, \lean{block_decomp_of_cycles}),
applied to its local graph (\lean{localGraph}), which joins each vertex of $K$ to its parents. Inside $K$ the children of
a vertex are its digon partners, so every vertex of $K$ has local degree at most three; a vertex outside $K$ has at most
one local neighbour, by~\ref{it:P2}; and the local graph has no triangle inside $K$, since such a triangle would consist
of digons and contain a transitive triangle. The cycle pieces and the final relative $2$-colouring then have private
outside parents.
\end{outline}

The blocks are usually far from minimum pieces; Theorem~PC is not a shortest-cycle argument. Its natural blocks are the
components of the digon graph.

\begin{corollary}[Theorem PC for digon components]\label{cor:PC-digon}
\notformalized
Let $N$ be loopless and TT-free. Suppose that
\begin{enumerate}[label=(H\arabic*),leftmargin=3em]
\item every vertex has at most one one-way parent, and
\item no vertex has one-way arcs to two vertices of the same component of the digon graph $G_2(N)$,
\end{enumerate}
and that every component $K$ of $G_2(N)$ has no one-way arc inside it and only vertices with at most three parents, or
induces a path of $G_2(N)$, or is such that $N[K]$ has an elimination certificate (for instance a cycle of $G_2(N)$ of length $m$ whose
one-way arcs inside $K$ all have the same offset $k$ with $2k\not\equiv\pm1\pmod m$, by \cref{prop:CC} with the mirror
sweep; TT-freeness gives $3\le k\le m-3$). Then $N$ satisfies \eqref{eq:star}.
\end{corollary}

\begin{informalproof}
Take the components of $G_2(N)$ as blocks. The digon partners of a vertex lie in its own component, so its parents
outside its component are one-way; by~(H1) there is at most one, which is~\ref{it:P1}. A vertex outside a component $K$
has only one-way arcs into $K$, at most one by~(H2), which is~\ref{it:P2}. For~\ref{it:P3}: a component without one-way
arcs inside it has every internal arc in a digon, and with the in-degree bound and TT-freeness it is of kind~(a). If $K$
induces a path $k_0,\dots,k_{L-1}$ of $G_2(N)$, its consecutive pairs are digons, and every other arc inside $K$ is
one-way, because a digon inside $K$ is an edge of $G_2(N)[K]$, which is the path. Each vertex is the head of at most one
one-way arc, by~(H1), and the tail of at most one one-way arc into $K$, by~(H2); so $K$ is a chorded digon path, of
kind~(b). The remaining components are of kind~(c) by hypothesis. Theorem~PC applies.
\end{informalproof}

Every chordless cycle of $G_2(N)$ is of kind~(a), and also of kind~(c) by \cref{lem:cycle-cert}\ref{it:cycle}; it has
length at least four, since $G_2(N)$ is triangle-free. Whether every TT-free chorded digon cycle is of kind~(c) is
\cref{conj:K}.

\paragraph{Why it matters.} Theorem~PC covers relations that none of Theorems~N3, U3 and~T covers, and puts no bound on
out-degrees: a vertex may send one one-way arc into each of arbitrarily many blocks. For~(T), side information built
from two-way exchanges inside blocks, linked by one-way arcs that are private in the sense of~\ref{it:P1}
and~\ref{it:P2}, gives no coding gain. With the digon components as blocks, a construction needs, on a large scale,
vertices that receive several one-way arcs or send several into one component, or components without certificates. The
relation $\Ex$ of \cref{rem:X15} has vertices of the first kind, and yet it has a global certificate.

\subsection{Cayley relations: Corollary CC}\label{sec:CC}

\begin{corollary}[Corollary CC: Cayley relations]\label{cor:CC}
\formal{Cayley.lean}{cayN, coset_steps, mem_coset_iff, rateBound_cayley, card_sq_le_cayley}
Let $G$ be a finite group and let $s,t\in G$ be such that $S=\{s,s^{-1},t\}$ has three elements and does not contain
$1$. Let $N$ be the Cayley relation with arcs $g\to gx$ for $x\in S$, so $N(q)=\{qs^{-1},qs,qt^{-1}\}$. If $N$ is
TT-free, then $N$ satisfies \eqref{eq:star}, and every storage code $C$ for $N$ over a nonempty finite alphabet
satisfies $\card C^2\le\card\alpha^{\card G}$.
\end{corollary}

\begin{outline}
The relation is loopless because $1\notin S$. If $t^2=1$, every arc is reciprocated, the underlying graph has degree at
most three, and Theorem~U3 applies. Otherwise take the left cosets $g\langle s\rangle$ as blocks. The order $m$ of $s$ is at least four: $m\ne1,2$ because $s\ne1$ and $s\ne s^{-1}$, and $m\ne3$
because $g\to gs\to gs^2=gs^{-1}$ with $g\to gs^{-1}$ would be a transitive triangle. The only parent of $q$ outside its
coset is $qt^{-1}$, and the only possible child outside it is $qt$; so~\ref{it:P1} and~\ref{it:P2} hold. Index the coset of
$g$ by $j\mapsto gs^j$ (\lean{mem_coset_iff}); in-neighbours through $s^{\pm1}$ then sit at neighbouring positions
(\lean{coset_steps}). If $t\notin\langle s\rangle$, its internal arcs join consecutive positions,
and \cref{lem:cycle-cert}\ref{it:cycle} applies. If $t=s^k$, its internal arcs are $gs^{j\pm1}\to gs^j$ and
$gs^{j-k}\to gs^j$. Here $k\ne0,\pm1$ because $\card S=3$ and $1\notin S$; $k\ne\pm2$ and $2k\not\equiv\pm1$ because
each would close a transitive triangle. So $3\le k\le m-3$ and \cref{prop:CC} applies, with the generator $s$ or, for the
mirror sweep, $s^{-1}$ and the offset $m-k$. Theorem~PC concludes.
\end{outline}

If instead $S$ has at most two elements and does not contain~$1$, every vertex has at most two parents, and Theorem~T
applies when $N$ is TT-free. For example,
$\Cay(\Z_9,\{1,-1,3\})$, with $N(q)=\{q-1,q+1,q-3\}$, is covered by Corollary~CC (one block, $m=9$,
$k=3$, $\delta=5$). It is covered by none of the earlier results: its underlying graph is $4$-regular, so neither
Theorem~N3, even with \cref{lem:fewer-arcs}, nor Theorem~U3 applies; it is strongly connected with every in-degree equal
to three, so Theorem~T and Corollary~$\Tint$ do not apply; and it is not $2$-dicolourable, since the two ends of a
digon must lie in different acyclic sets and its digons form a $9$-cycle. It also shows that the private pieces of Theorem~N3 can fail to exist (\cref{prop:no-private}).

\paragraph{Why it matters.} Cayley relations are the side information of uniform, translation-invariant schedules. For
every finite group, the pattern ``two-way exchange along one generator, one-way transfer along another'' gives no
coding gain at any size. For~(T), symmetric constructions of this shape are excluded, whatever the group.

\section{Barrier bounds, guessing numbers and index codes}\label{sec:consequences}

The relative form of Theorem~D gives exact values on triangle-free graphs of maximum degree three: the guessing number
is the fractional matching number, and the broadcast rate of index coding is the number of vertices minus it. We end the
section with a general consequence for relations that become one of our classes after deleting a few vertices.

For a graph $G$ and $S\subseteq V$ let $I(S)=\{v\notin S:N_G(v)\subseteq S\}$, the set of isolated vertices of $G-S$
(\lean{isolatedOut}).

\begin{theorem}[Barrier bound]\label{thm:barrier}
\formal{Barrier.lean}{isolatedOut, barrier_ineq, log_card_le_barrier}
Let $G$ be a triangle-free graph of maximum degree at most~$3$ and $S\subseteq V$. Every polymatroid $h$ valid for $N_G$
satisfies
\[
  2h(V)\;\le\;2\sum_{v\in S}h(v)+\sum_{v\notin S\cup I(S)}h(v),
\]
and every storage code $C$ for $G$ over a nonempty finite alphabet $\alpha$ satisfies
$2\log\card C\le(\card V+\card S-\card{I(S)})\log\card\alpha$.
\end{theorem}

\begin{outline}
Let $U=S\cup I(S)$. Theorem~N3 gives a decomposition from the context $U$ (\lean{n3_decomp}), and the relative form of
Theorem~D gives $2(h(V)-h(U))\le\sum_{v\notin U}(h(U\cup\{v\})-h(U))\le\sum_{v\notin U}h(v)$. The vertices of $I(S)$
have all their neighbours in $S$, so $h(U)=h(S)\le\sum_{v\in S}h(v)$ by the closure lemma and subadditivity. For codes,
apply this to the entropy of a uniformly random codeword, as in \cref{thm:codes}.
\end{outline}

\begin{definition}[Fractional matchings]\label{def:fracmatching}
\formal{Matching.lean}{IsFracMatching, fracWeight, fracMatchingNumber}
A \emph{fractional matching} of a graph $G$ is a symmetric function $y\colon V\times V\to\R_{\ge0}$ that vanishes on
non-adjacent pairs and satisfies $\sum_uy(v,u)\le1$ for every $v$. Its \emph{weight} is $\frac12\sum_{u,v}y(u,v)$, each
edge counted once. The \emph{fractional matching number} $\nustar(G)$ is the supremum of the weights.
\end{definition}

\begin{proposition}[The fractional Tutte--Berge formula]\label{prop:fracTB}
\formal{Matching.lean}{fracWeight_le_barrier, exists_doubleCover_matching, dcFrac, dcFrac_isFracMatching, fracWeight_dcFrac, fracMatchingNumber_spec, fracMatchingNumber_eq_min_barrier}
For every finite simple graph $G$, the supremum $\nustar(G)$ is attained by a half-integral fractional matching, and
\[
  2\nustar(G)\;=\;\min_{S\subseteq V}\bigl(\card V+\card S-\card{I(S)}\bigr).
\]
\end{proposition}

The formula is classical (see~\cite{ScheinermanUllman1997}); we formalize it because the next theorem needs both
directions.

\begin{outline}
Weak duality (\lean{fracWeight_le_barrier}): the vertices of $I(S)$ send all their weight into $S$, which receives at most
$\card S$, and every other vertex carries weight at most one, so twice the weight is at most
$\card S+(\card V-\card{I(S)})$. For the converse, take a set $X_0$ of largest deficiency $d=\card{X_0}-\card{N(X_0)}$,
where $N(X)$ is the set of neighbours of $X$. Hall's theorem~\cite{Hall1935}, in its Mathlib form, applied to the
bipartite double cover of $G$ with $d$ dummy vertices added on one side, gives a set $D\subseteq V$ and an injective map
$g$ on $D$ sending each vertex to a neighbour, with at most $d$ vertices outside $D$
(\lean{exists_doubleCover_matching}). Putting weight $\frac12$ on each pair $\{v,g(v)\}$, once per direction, gives a
half-integral fractional matching of weight $\card D/2$ (\lean{dcFrac}, \lean{fracWeight_dcFrac}). For
$S=N(X_0)\setminus X_0$ the set $I(S)$ contains $X_0\setminus N(X_0)$, which gives
$\card D\ge\card V+\card S-\card{I(S)}$. With weak duality, both bounds are equalities (\lean{fracMatchingNumber_spec}).
\end{outline}

\begin{theorem}[The guessing number is the fractional matching number]\label{thm:gn}
\formal{Matching.lean}{log_card_le_fracMatchingNumber, exists_code_fracMatchingNumber, guessingNumber_eq_fracMatching}
\begin{enumerate}[label=(\alph*)]
\item\label{it:gn-lower} For every finite simple graph $G$ and every nonempty finite set $\beta$, some storage code for
$G$ over the alphabet $\beta\times\beta$ has at least $\card{\beta\times\beta}^{\nustar(G)}$ words.
\item\label{it:gn-upper} If $G$ is triangle-free of maximum degree at most~$3$, every storage code for $G$ over every
nonempty finite alphabet $\alpha$ has $\card C\le\card\alpha^{\nustar(G)}$, and over every alphabet $\beta\times\beta$
some storage code has exactly $\card{\beta\times\beta}^{\nustar(G)}$ words.
\end{enumerate}
\end{theorem}

The Lean statements are logarithmic, for instance $\nustar(G)\log\card{\beta\times\beta}\le\log\card C$, with the
convention $\log0=0$. They give the counts above whenever $\nustar(G)>0$ and $\card\beta\ge2$; in the remaining cases the
claims about $\beta\times\beta$ ask only for one word, and a code with a single word is a storage code.

\begin{outline}
\ref{it:gn-lower} Take $D$ and $g$ as in \cref{prop:fracTB}. The code consists of the words $x$ in which, for every
$v\in D$, the second half of $x_v$ equals the first half of $x_{g(v)}$, and all other halves are a fixed symbol: each
vertex of $D$ shares one half of its symbol with its neighbour $g(v)$, and each vertex of $g(D)$ receives one half from
its preimage. Every symbol is then determined by the neighbours, and the $\card\beta^{\card D}=
\card{\beta\times\beta}^{\nustar}$ choices of the shared halves give distinct words
(\lean{exists_code_fracMatchingNumber}).
\ref{it:gn-upper} \Cref{thm:barrier} with a set $S$ attaining the minimum of \cref{prop:fracTB}
(\lean{log_card_le_fracMatchingNumber}); equality by~\ref{it:gn-lower}.
\end{outline}

\begin{corollary}[Guessing numbers]\label{cor:gn}
\formal{Rates.lean}{guessingNumberAt, guessingNumber, isStorageCode_map_equiv, isStorageCode_map_injective, logb_card_le_fracMatchingNumber, guessingNumberAt_le, guessingNumberAt_sq, guessingNumberAt_ge, guessingNumber_eq, tendsto_guessingNumberAt}
For a finite simple graph $G$ and $q\ge2$, let $\gn(G,q)$ be the supremum of $\log_q\card C$ over the storage codes $C$
for $G$ over the alphabet $\{0,\dots,q-1\}$, and let $\gn(G)=\sup_{q\ge2}\gn(G,q)$. If $G$ is triangle-free of maximum
degree at most~$3$, then
\begin{enumerate}[label=(\alph*)]
\item\label{it:gnq} $\gn(G,q)\le\nustar(G)$ for every $q\ge2$;
\item\label{it:gnsq} $\gn(G,r^2)=\nustar(G)$ for every $r\ge2$;
\item\label{it:gnsup} $\gn(G)=\nustar(G)$;
\item\label{it:gnlim} $\gn(G,q)\to\nustar(G)$ as $q\to\infty$.
\end{enumerate}
\end{corollary}

\begin{outline}
The Lean definitions (\lean{guessingNumberAt}, \lean{guessingNumber}) take the supremum with the convention
$\log0=0$ for the empty code, which does not change it, since one-word codes give $0$. \ref{it:gnq} Divide the
logarithmic bound of \cref{thm:gn}\ref{it:gn-upper} by $\log q>0$ (\lean{logb_card_le_fracMatchingNumber},
\lean{guessingNumberAt_le}). \ref{it:gnsq} For $\card\beta=r$, the code of \cref{thm:gn}\ref{it:gn-lower} over
$\beta\times\beta$ has at least $(r^2)^{\nustar(G)}$ words, and a bijection $\beta\times\beta\to\{0,\dots,r^2-1\}$
carries it to a storage code over that alphabet (\lean{isStorageCode_map_equiv}); with~\ref{it:gnq},
$\gn(G,r^2)=\nustar(G)$ (\lean{guessingNumberAt_sq}). \ref{it:gnsup} Every $\gn(G,q)$ is at most $\nustar(G)$, and
$\gn(G,4)=\nustar(G)$ (\lean{guessingNumber_eq}). \ref{it:gnlim} For $q\ge4$ let $r=\lfloor\sqrt q\rfloor\ge2$. An
injection $\{0,\dots,r-1\}^2\to\{0,\dots,q-1\}$ carries the code of \cref{thm:gn}\ref{it:gn-lower} over
$\{0,\dots,r-1\}^2$ to a storage code over $\{0,\dots,q-1\}$ (\lean{isStorageCode_map_injective}), so
$\gn(G,q)\ge\nustar(G)\log(r^2)/\log q$ (\lean{guessingNumberAt_ge}). Since $q<(r+1)^2\le4r^2$, this is at least
$\nustar(G)(1-\log4/\log q)$, using $\nustar(G)\ge0$ (\lean{fracMatchingNumber_nonneg}); with the upper bound
of~\ref{it:gnq}, $\gn(G,q)\to\nustar(G)$ (\lean{tendsto_guessingNumberAt}).
\end{outline}

The restriction to squares in~\ref{it:gnsq} cannot be dropped: by a theorem of Atkins, Rombach and
Skerman~\cite{AtkinsRombachSkerman2017}, for every $q\ge2$ that is not a perfect square and every sufficiently large odd
$n$, the odd cycle $C_n$, which lies in this class and has $\nustar(C_n)=n/2$, has $\gn(C_n,q)=(n-1)/2+\log_qa<n/2$, where
$a$ is the largest divisor of $q$ that is at most $\sqrt q$.

The code of part~\ref{it:gn-lower} realizes Christofides and Markström's fractional clique cover
strategy~\cite{ChristofidesMarkstrom2011}, restated in storage language by Mazumdar, McGregor and
Vorotnikova~\cite{MMV2019}: on a triangle-free graph the cliques are the vertices and the edges, and the fractional clique
cover number is $\card V-\nustar(G)$ (\cref{cor:broadcast}\ref{it:bc-cover}). Christofides and Markström conjectured
that this strategy is optimal on every undirected graph (see~\cite[Conjecture~1]{Dang2015}) and proved it for perfect
graphs and odd cycles~\cite{ChristofidesMarkstrom2011}; Mazumdar, McGregor and Vorotnikova proved it for prisms and
certain chorded cycles~\cite{MMV2019}. The conjecture is false in general: Baber, Christofides, Dang, Vaughan and Riis disproved
it~\cite{Baber2017} (Dang's thesis gives a counterexample on ten vertices~\cite[Section~4.1]{Dang2015}), and Cameron, Dang
and Riis~\cite{CameronDangRiis2016} found triangle-free counterexamples, such as the Higman--Sims graph. For the Clebsch
graph, of degree five, they proved $10\le\gn\le11$, and the fourth paper determined $\gn=10>8=\nustar$
\cite[Corollary~6.11]{FourthPaper2026}. \Cref{thm:gn,cor:gn} prove the conjecture for triangle-free graphs of maximum
degree three, whether the asymptotic guessing number is taken as the supremum over $q$, as here, or as the limit as
$q\to\infty$, in which form the conjecture is often stated (\cref{cor:gn}\ref{it:gnlim}).

\begin{definition}[Index codes]\label{def:indexcode}
\formal{IndexCoding.lean}{IndexCode}
An \emph{index code} for $N$ over an alphabet $\alpha$ with a finite message set $M$ is an encoding map $E\colon\alpha^V\to M$
together with decoders $D_v$, $v\in V$, such that $D_v(E(x),x)=x_v$ for every word $x$, where $D_v(m,x)$ depends only on
$m$ and on the restriction of $x$ to $N(v)$.
\end{definition}

\begin{theorem}[Index codes and fractional matchings]\label{thm:index}
\formal{IndexCoding.lean}{IndexCode.isStorageCode_fibre, log_card_message_ge, exists_indexCode}
\begin{enumerate}[label=(\alph*)]
\item If $G$ is triangle-free of maximum degree at most~$3$, every index code for $N_G$ over a nonempty finite alphabet
$\alpha$ has $\log\card M\ge(\card V-\nustar(G))\log\card\alpha$.
\item For every finite simple graph $G$ and every finite abelian group $\beta$, there is an index code for $N_G$ over
$\beta\times\beta$ with $\log\card M=(\card V-\nustar(G))\log\card{\beta\times\beta}$.
\end{enumerate}
\end{theorem}

\begin{outline}
(a) The words with a common message form a storage code (\lean{IndexCode.isStorageCode_fibre}), as in the fourth
paper~\cite[Remark~4.9(b)]{FourthPaper2026}; each has at most $\card\alpha^{\nustar}$ words by \cref{thm:gn}, and the
$\card M$ fibres cover $\alpha^V$. (b) With $D$ and $g$ as in \cref{prop:fracTB}, broadcast, for every $v$, the
difference between the second half of $x_v$ and the first half of $x_{g(v)}$ when $v\in D$, and the second half of $x_v$
otherwise, and broadcast the first halves of the vertices outside $g(D)$. A vertex recovers the second half of its symbol
from its message and the first half of $x_{g(v)}$, and, if $v=g(u)$, the first half from the message of $u$ and the second
half of $x_u$; both $g(v)$ and $u$ are neighbours. The message set has $\card\beta^{2\card V-\card D}$ elements.
\end{outline}

The duality between storage codes and index codes behind part~(a) is that of Mazumdar~\cite{Mazumdar2014}, Shanmugam and
Dimakis~\cite{ShanmugamDimakis2014}, and Arbabjolfaei and Kim~\cite{ArbabjolfaeiKim2015}.

\begin{corollary}[The broadcast rate and fractional clique covers]\label{cor:broadcast}
\formal{Rates.lean}{IndexCode.ofInjective, IndexCode.restrict, broadcastBits, broadcastBits_nonempty, broadcastBits_ge, broadcastBits_le_succ, broadcastBits_even_le, broadcastBits_even, fracMatchingNumber_nonneg, fracMatchingNumber_le_card, tendsto_broadcastBits, IsFracCover, coverWeight, fracCliqueCoverNumber, fracCliqueCoverNumber_eq}
Let $G$ be a finite simple graph. For $t\ge0$ let $\beta_t$ be the least $n$ such that $N_G$ has an index code over the
alphabet $\{0,1\}^t$ with message set $\{0,1\}^n$.
\begin{enumerate}[label=(\alph*)]
\item\label{it:bc-every} For all $t,s\ge0$, $\beta_t\le\beta_{t+1}$ and $\beta_{2s}\le2(\card V-\nustar(G))s$.
\item\label{it:bc-rate} If $G$ is triangle-free of maximum degree at most~$3$, then $(\card V-\nustar(G))t\le\beta_t$ for
every $t$, $\beta_{2s}=2(\card V-\nustar(G))s$ for every $s$, and $\beta_t/t\to\card V-\nustar(G)$ as $t\to\infty$.
\item\label{it:bc-cover} Call weights $y_v\ge0$ on the vertices and $y(u,v)=y(v,u)\ge0$ on pairs of vertices, vanishing on
non-adjacent pairs, a \emph{fractional cover by vertices and edges} if $y_v+\sum_uy(v,u)\ge1$ for every $v$, and let its
weight be $\sum_vy_v+\frac12\sum_{u,v}y(u,v)$, each edge counted once. The infimum of these weights is
$\card V-\nustar(G)$.
\end{enumerate}
\end{corollary}

\begin{outline}
\ref{it:bc-every} The identity code shows that $\beta_t$ exists (\lean{broadcastBits_nonempty}). An index code over
$\{0,1\}^{t+1}$ restricts to $\{0,1\}^t$: append a fixed bit to every symbol, and drop it after decoding
(\lean{IndexCode.restrict}, \lean{broadcastBits_le_succ}). \Cref{thm:index}(b) with $\beta=\F_2^s$ gives an index code
over $\beta\times\beta$, which is $\{0,1\}^{2s}$ up to a bijection, whose message set has $2^n$ elements with
$n=2(\card V-\nustar(G))s$; a bijection of message sets (\lean{IndexCode.ofInjective}) gives
\lean{broadcastBits_even_le}. \ref{it:bc-rate} \Cref{thm:index}(a) with $\alpha=\{0,1\}^t$ gives the lower bound
(\lean{broadcastBits_ge}), and with~\ref{it:bc-every} the equality for even $t$ (\lean{broadcastBits_even}). For odd $t$,
$\beta_t\le\beta_{t+1}=(\card V-\nustar(G))(t+1)$, so $\card V-\nustar(G)\le\beta_t/t\le(\card V-\nustar(G))(1+1/t)$
for all $t\ge1$, where $0\le\nustar(G)\le\card V$ (\lean{fracMatchingNumber_nonneg}, \lean{fracMatchingNumber_le_card},
\lean{tendsto_broadcastBits}). \ref{it:bc-cover} In Lean the covers are \lean{IsFracCover}, their weight is
\lean{coverWeight}, the infimum is \lean{fracCliqueCoverNumber}, and the claim is \lean{fracCliqueCoverNumber_eq}. Take $S$ attaining the minimum in
\cref{prop:fracTB}, and put $z_v=1$ on $I(S)$, $z_v=0$ on $S$ and $z_v=\frac12$ elsewhere. An edge at a vertex of $I(S)$
ends in $S$, so $z_u+z_v\le1$ on every edge, and every cover has weight at least
$\sum_vz_v(y_v+\sum_uy(v,u))\ge\sum_vz_v=\frac12(\card V-\card S+\card{I(S)})=\card V-\nustar(G)$. Conversely, an optimal
fractional matching $x$ gives the cover with $y(u,v)=x(u,v)$ and $y_v=1-\sum_ux(v,u)$, of weight $\card V-\nustar(G)$.
\end{outline}

The limit in~\ref{it:bc-rate} is the broadcast rate of Blasiak, Kleinberg and Lubetzky~\cite{BKL2013}, so the broadcast
rate of a triangle-free graph of maximum degree three is $\card V-\nustar(G)$. In a triangle-free graph the cliques are
the vertices and the edges, so the infimum in~\ref{it:bc-cover} is the fractional clique cover number, which therefore
equals $\card V-\nustar(G)$. The bound of \cref{thm:index}(b) is thus the fractional clique cover upper bound on the
broadcast rate~\cite{BKL2013}, realized over $\beta\times\beta$ by a double-cover matching. For some graphs of the class
these values were known: Blasiak, Kleinberg and Lubetzky determined the broadcast rate of cycles and of the cubic
circulants, and computed that of the Petersen graph numerically~\cite{BKL2013}.

\paragraph{Why it matters.} For triangle-free subcubic side information, \cref{thm:barrier,thm:gn,thm:index} and
\cref{cor:gn,cor:broadcast} determine the guessing number and the broadcast rate; the bounds hold for all codes and all
alphabets and are attained over the alphabets $\beta\times\beta$. In the language of the fourth paper's crux
(\cref{sec:construction}), on such side information an index code must broadcast at least $\card V-\nustar(G)\ge\card
V/2$ symbols, and the extremal codes are fractional pairings: there is no room for a coding gain.

\subsection{Relations near a class}\label{sec:paid}

All our classes are defined by exact conditions. A few exceptional vertices cost only what they could carry.

\begin{corollary}[Paid deletion]\label{cor:paid}
\formal{PairedDistance.lean}{card_sq_le_of_paidDeletion}
Let $N$ be a recovery relation on $V$ and $Z\subseteq V$. If $N-Z$ satisfies \eqref{eq:star}, then every storage code
$C$ for $N$ over a nonempty finite alphabet $\alpha$ satisfies $\card C^2\le\card\alpha^{\card V+\card Z}$.
\end{corollary}

\begin{outline}
In the formal statement, $N-Z$ is the relation \lean{DirectedRecovery.inducedNeighborhood}~$N\,(V\setminus Z)$ induced
on $V\setminus Z$, as in \emph{Directed Recovery Structures}. By \cref{thm:codes}
(\lean{MotionDepthCodes.card_sq_le_of_rateBound_half}), every storage code for $N-Z$ has at most
$\card\alpha^{(\card V-\card Z)/2}$ words. The paid-conditioning lemma of
\emph{Directed Recovery Structures}~\cite[Lemma~2.2]{DirectedRecovery2026}
(\lean{DirectedRecovery.storageCode_conditioning_bound}) fixes the values on $Z$: every fibre is a storage code for
$N-Z$, and there are at most $\card\alpha^{\card Z}$ fibres, so $\card C\le\card\alpha^{\card Z+(\card V-\card Z)/2}$.
Squaring gives the claim.
\end{outline}

For $\card\alpha\ge2$ the conclusion says that $\log_{\card\alpha}\card C\le(\card V+\card Z)/2$.

Combined with \cref{lem:tt-paired}, the corollary shows that relations with
few short paired paths are as good as TT-free ones: if deleting the loops, the middle vertices of transitive triangles and
a further set $Z'$ leaves a relation in one of our classes, then storage codes have rate at most
$\frac12+(b_E(3)+\card{Z'})/(2\card V)$.

\paragraph{Why it matters.} This is how the exact theorems of this paper apply to the asymptotic setting of the
paired-distance question~\cite[Question~9.1]{DirectedRecovery2026}: a construction must keep its cone relations away
from every class here by a linear number of vertices (\cref{sec:construction}).

\section{Limits of the method}\label{sec:limits}

The decompositions above are produced by rules: contract a shortest cycle (Theorems~N3 and~U3), a shortest directed cycle
of in-degree-two vertices (Theorem~T), or a whole block (Theorem~PC). This section shows where such rules stop.
\Cref{prop:no-private} and \cref{ex:D9} are formalized, by kernel evaluation. The remarks report computations,
labelled as numerical, a cited formula and a heuristic; where parts of them are Lean-verified, they say so, and none of
them is used in a proof. The computations are reproduced by the script \leanfile{tools/verify_remarks.py}, described in
\leanfile{tools/README.md}.

\subsection{Private pieces need not exist}\label{sec:no-private}

The induction of Theorem~N3 removes a piece whose outside parents are private. In directed relations of degree three
such pieces need not exist, even when the half bound holds.

\begin{proposition}[No private proper pieces]\label{prop:no-private}
\formal{Limits.lean}{cay9, no_private_piece}
Let $N$ be the relation of $\Cay(\Z_9,\{1,-1,3\})$, so $N(q)=\{q-1,q+1,q-3\}$ and the children of $q$ are $q+1$, $q-1$
and $q+3$. There is no set $H$ with $\emptyset\ne H\ne\Z_9$ such that
\begin{enumerate}[label=(\alph*)]
\item every vertex of $H$ has at most one parent outside $H$, and
\item every vertex outside $H$ is a parent of at most one vertex of $H$.
\end{enumerate}
\end{proposition}

\begin{outline}
The formal proof evaluates the statement for all $512$ subsets $H$ of $\Z_9$ in the kernel (\texttt{decide +kernel}).
The following argument explains the result. Suppose $H$ exists. Since $\emptyset\ne H\ne\Z_9$, some $a\in H$ has
$a-1\notin H$. By~(a) at $a$, two of its parents $a-1,a+1,a-3$ lie in $H$, so $a+1\in H$ and $a-3\in H$. The children of
$a-1\notin H$ are $a-2$, $a$ and $a+2$; since $a\in H$, (b) gives $a-2\notin H$ and $a+2\notin H$. Now the parents of
$a+1$ are $a$, $a+2$ and $a-2$, and only $a$ lies in $H$, contradicting~(a). The residues $a-3,\dots,a+2$ are distinct
modulo~$9$, so the argument is valid.
\end{outline}

Every piece of a decomposition in context $\emptyset$ satisfies~(a), and a piece whose outside parents are private
satisfies~(b). So on this relation a decomposition must start with the whole vertex set or with a piece whose outside
parents are shared. Both happen. Corollary~CC gives a single piece; explicitly, the records $\{0,4\}$, $\{1,3,6,8\}$,
$\{2,5,7\}$ with the list $0,1,2,3,4,7,5,6,8$ form a global elimination certificate (\lean{z9_elimOn}), which gives
\eqref{eq:star} directly (\lean{rateBound_of_elimOn_univ}, \lean{rateBound_cay9}). And nested charges give a proper
first piece, as follows.

\begin{example}[A proper piece with shared outside parents]\label{ex:D9}
\formal{Limits.lean}{exH1, ex72Piece, ex72Piece4, ex72Piece78, ex72_decomp}
In the relation of \cref{prop:no-private}, let $H_1=\{0,1,2,3,5,6\}$ with records $X=\{0,2,6\}$, $Y=\{1,3,5\}$ and an
empty record $Z$, and the list $2,3,5,6,0,1$ with supplying records $Y,X,X,Z,Y,Z$. The outside parents are $8$ for the
vertices $2$ and $0$, $4$ for $3$ and $5$, and $7$ for $6$ and $1$; each outside vertex is shared by two steps with the
same supplying record, so the certificate is co-supplied (\cref{lem:cosupply}). For instance, at the step of $5$ the
revealed set is $\{2,3\}$ and $N(5)=\{4,6,2\}$ lies in $\{2,3\}\cup X\cup\{4\}$. The relation that remains,
on $\{4,7,8\}$ with $H_1$ as context, is decomposed by the pieces $\{4\}$ and $\{7,8\}$, with records $\{4\},\emptyset$
and $\{7\},\{8\}$. So $H_1,\{4\},\{7,8\}$ is a decomposition of $N$.
\end{example}

\begin{outline}
The three piece certificates are written out (\lean{ex72Piece}, \lean{ex72Piece4}, \lean{ex72Piece78}); their
conditions, including co-supply (\lean{relCharged_of_cosupply}), are checked by kernel evaluation (\texttt{decide}), and
\lean{ex72_decomp} assembles them.
\end{outline}

\paragraph{Why it matters.} For an impossibility proof, \cref{prop:no-private} shows that the private-piece induction of
Theorem~N3 cannot prove half bounds on directed side information of degree three, even on a single circulant: an argument
must use whole blocks, as in Theorem~PC, or outside parents shared by several steps, which is what the nested charging of
Lemma~E provides.

\subsection{Minimum nonempty cores are not enough}\label{sec:cores}

Call $H\subseteq V$ a \emph{core} if every vertex of $H$ has at most one parent outside $H$. A piece in context
$\emptyset$ is a core, and in a cubic graph the minimum nonempty cores are the shortest cycles, which is why Theorem~N3
can always contract a minimum nonempty core. In directed relations this fails.

\begin{remark}[Numerical: a $15$-vertex relation whose minimum nonempty core is not a piece]\label{rem:X15}
Let $\Ex$ be the relation (\lean{x15}) on $\{0,\dots,14\}$ with
\begin{center}
\small
\begin{tabular}{@{}l*{8}{c}@{}}
\toprule
$q$ & 0 & 1 & 2 & 3 & 4 & 5 & 6 & 7\\
$N(q)$ & $\{5,1,6\}$ & $\{0,2,7\}$ & $\{1,3,8\}$ & $\{2,4,6\}$ & $\{3,5,7\}$ & $\{4,0,8\}$ & $\{14,0,10\}$ & $\{13,4,9\}$\\
\midrule
$q$ & 8 & 9 & 10 & 11 & 12 & 13 & 14 & \\
$N(q)$ & $\{9,14,5\}$ & $\{6,11,7\}$ & $\{9,12\}$ & $\{2,13,12\}$ & $\{14,11\}$ & $\{1,8,11\}$ & $\{3,13,12\}$ & \\
\bottomrule
\end{tabular}
\end{center}
The relation $\Ex$ is loopless and TT-free, and every vertex has at most three parents and at most three children
(\lean{x15_props}). Every nonempty core has at least six vertices, and the only core with six vertices is
$C=\{0,\dots,5\}$ (\lean{x15_cores}). The set $C$ is a core, its cyclically consecutive vertices form digons, and its
outside parents $6$, $7$, $8$ each have two antipodal children: $6\to0,3$, $7\to1,4$ and $8\to2,5$ (\lean{x15_hexagon}).
So $C$ is the unique minimum nonempty core, and, as every vertex has at most three parents, it is a bidirected
$6$-cycle. A computation (\leanfile{tools/verify_remarks.py}) shows that $\Ex$ is strongly connected and not
$2$-dicolourable, and an exhaustive search over the partitions of $C$ into at most seven records, all elimination orders
and all supplying records finds no piece certificate for $C$ in context $\emptyset$, even when the steps charged to an
outside parent are only required to be nested (co-supplied outside parents are a special case). More records cannot
help: $C$ has six vertices, so a certificate with more than seven records has at least two empty records, which can be
merged. So $C$ is not a piece in context $\emptyset$, with any number of records. Nevertheless $\Ex$ has a
decomposition, with a single piece: a global elimination certificate with three records, obtained by giving the vertices
$0,\dots,14$ the records $0,1,0,1,0,1,1,1,0,0,0,1,0,0,2$ and using the list $0,4,5,7,1,9,8,2,11,12,10,6,3,13,14$
(\lean{x15_elimOn}). So \eqref{eq:star} holds for $\Ex$ (\lean{rateBound_x15}), but no rule that starts by contracting a
minimum nonempty core can prove it. The claims that $\Ex$ is strongly connected and not $2$-dicolourable and that $C$ is
not a piece are numerical; the other claims of this remark are Lean-verified as cited.
\end{remark}

In the undirected setting the antipodal pattern cannot occur on a shortest cycle: an outside vertex adjacent to $c_0$ and
$c_3$ on a $6$-cycle closes the $5$-cycle $w,c_0,c_1,c_2,c_3$. In $\Ex$ the arc $6\to3$ is one-way
($0\leftrightarrow6$ is a digon); the $5$-cycle $6,0,1,2,3$ of the underlying graph is not a core, since vertex~$6$ has
the two parents $10$ and $14$ outside it, and no nonempty core has fewer than six vertices (\lean{x15_cores}).
The relation lies outside Theorems~U3 and~T: vertex $6$ has five neighbours in the underlying graph, and the vertices with
three parents include the digons of $C$. Its digon components violate~(H1) of \cref{cor:PC-digon}, since vertex $6$ has
the two one-way parents $10$ and $14$ (\lean{x15_oneWay_parents}); Theorem~PC applies with the whole vertex set as one
block, through the global certificate.

\subsection{Sweeps and global certificates}\label{sec:sweeps}

For a chorded digon cycle $K=\{k_0,\dots,k_{L-1}\}$ (indices modulo $L$) such that $N[K]$ is loopless and TT-free, the
\emph{sweep} from $k_z$ in direction $d=\pm1$ is the elimination list $k_z,k_{z+d},\dots,k_{z+(L-1)d}$ with the finest
records that make the later parents of each vertex lie in one record: the classes of the equivalence relation that
merges these later parents, for every vertex. The sweep is \emph{valid} if no vertex lies in the class of its later
parents; since $N[K]$ is loopless, it is then an elimination certificate. The proof of \cref{prop:CC}\ref{it:mirror}
shows that such a cycle (with $N[K]$ loopless and TT-free) whose chords all have one offset $k$ with
$2k\not\equiv\pm1\pmod L$ has a valid sweep, along the cycle or, for the mirror sweep, in the opposite direction; this includes every loopless TT-free
circulant $\Cay(\Z_L,\{1,-1,k\})$.

\begin{remark}[Numerical: chorded cycles without a valid sweep]\label{rem:sweeps}
An adversarial local search over chord sets found twelve TT-free chorded digon cycles, of lengths $13$, $15$ and $16$,
none of whose $2L$ sweeps is valid (\leanfile{tools/verify_remarks.py}). One of them has $L=13$ and the chords
\[
 0\to8,\ 1\to6,\ 3\to11,\ 4\to0,\ 5\to10,\ 6\to9,\ 7\to3,\ 8\to4,\ 9\to12,\ 10\to1,\ 11\to7,\ 12\to2.
\]
All twelve have elimination certificates with three records. For the cycle displayed, give the vertices
$0,\dots,12$ the records $0,2,0,1,0,2,1,2,1,2,1,2,1$; the list $11,9,7,6,5,4,10,1,2,0,8,3,12$ is then an elimination
certificate, which can be checked by hand. So these cycles are of kind~(c) in Theorem~PC; they only show that sweeps, the
method of \cref{prop:CC}, do not cover all chorded cycles.

All twelve cycles are Lean-verified to be loopless TT-free chorded digon cycles: on $\Z_L$, every vertex $v$ has the
parents $v\pm1$, at most three parents and at most three children (\lean{IsChordedCycle}), so beyond the digons of the
cycle there is at most one chord into and at most one chord out of every vertex. Their global elimination
certificates, which the Lean proofs construct with three records, and hence \eqref{eq:star}, are Lean-verified as well.
The displayed cycle is \lean{c13}: \lean{c13_spec} shows that the parents of each vertex are exactly its two cycle
neighbours and the tails of the displayed chords into it, and \lean{c13_isChordedCycle}, \lean{c13_elimOn} and
\lean{rateBound_c13} give the rest. The other eleven are \lean{cyc1}, \dots, \texttt{cyc11}, each with its
certificate \texttt{cyc$K$\_elimOn}; \lean{twelve_cycles} collects the twelve statements \lean{IsChordedCycle} and
\eqref{eq:star}. Apart from \lean{cyc1}, the declarations of these eleven cycles are not listed in the index
individually; their axioms are checked through \lean{twelve_cycles}. Only the claim that none of the $2L$ sweeps is
valid is numerical.
\end{remark}

\begin{remark}[Numerical: the Petersen graph has no global certificate]\label{rem:petersen}
Let $N$ be the neighbourhood relation of the Petersen graph. In a global elimination certificate of $N$
(\cref{def:elimcert}), written as an elimination list, the later neighbours of each vertex lie in one record, different
from the record of the vertex; when the number of records is not bounded, an elimination order admits a certificate
exactly when these forced equalities and inequalities of records are consistent, since the classes of the forced
equalities can then serve as the records. An exhaustive search over all elimination orders of $N$, which keeps the
forced equalities in a union--find structure and abandons an order as soon as a forced inequality fails, finds no
consistent order (\leanfile{tools/verify_remarks.py}). So $N$ has no global elimination certificate, with any number of
records. The Petersen graph, \lean{MotionDepthCodes.petersenGraph} of the fourth paper's library, is triangle-free
(\lean{petersen_cliqueFree}) and cubic, so Theorem~N3, by contrast, gives it a decomposition (\lean{petersen_decomp}).
Following the proof of Theorem~N3, the first piece is a $5$-cycle, whose outside parents (the spokes) are private, and the
second is the remaining $5$-cycle; this description of the pieces is not part of the formal statement. So
decompositions are strictly stronger than global certificates, and Theorem~D strictly stronger than
\cref{prop:global}.
\end{remark}

\subsection{Odd girth and degree}\label{sec:foldcubes}

\begin{remark}[Cited: folded cubes]\label{rem:foldcube}
The folded $(2k+1)$-cube $FQ_{2k+1}=\Cay(\F_2^{2k},\{e_1,\dots,e_{2k},\mathbf 1\})$, where $\mathbf1$ is the all-ones
vector, has $N=2^{2k}$ vertices, degree $2k+1$ and odd girth $2k+1$; for $k=2$ it is the Clebsch graph. Over $\F_2$, the
matrix $I+A$ has rank $\frac12(2^{2k}-2^k)$: this is the case of even dimension $2k$ of a proposition of Barg,
Schwartz and Yohananov~\cite[Proposition~13]{BSY2024}, whose Cayley graph also has the zero vector as a generator and
therefore has adjacency matrix $I+A$. Its kernel is a binary linear storage code of dimension
$N/2+2^{k-1}$; for $k=2,3$ this is also checked numerically (\leanfile{tools/verify_remarks.py}). For $k=2$ it is
Lean-verified in the fourth paper~\cite[Corollary~6.11]{FourthPaper2026}: its Clebsch graph is the folded $5$-cube
(\lean{MotionDepthCodes.clebschGraph_adj_iff}), and $\ker(I+A)$ has dimension $10=N/2+2$
(\lean{MotionDepthCodes.finrank_ker_clebsch}). So \eqref{eq:star} fails on the folded cubes, whose odd girth equals their
degree, for every $k\ge2$; it holds on other graphs whose odd girth equals their degree, such as the $5$-regular
Andrásfai graph $\Andr(5)=K_{14/5}$, of odd girth five (\cref{thm:circular}). In the symmetric case $\dist_E(L_p,R_p)$ is
the length of the shortest odd closed walk through $p$, so here $b_E(r)=0$ for $r<2k+1$. Hence, for every odd $r\ge3$, an exact half bound that assumes only $b_E(r)=0$ can hold at most for maximum degree
$r+1$: the folded $(r+2)$-cube violates it. For $r=3$ it fails already at maximum degree four (\cref{ex:deg4}), so the
degree bound of Theorem~N3 cannot be raised. The excess $2^{k-1}=\sqrt N/2$ is $o(N)$, which is consistent with the
asymptotic form of the paired-distance question~\cite[Question~9.1]{DirectedRecovery2026}.
\end{remark}

\begin{remark}[Heuristic: degree four]\label{rem:degree4}
In the undirected setting, the method needs each step to have at most one outside parent. If a vertex $v$ of a shortest
cycle has two outside neighbours $p$ and $p'$, the step that eliminates it costs $I(v;\{p,p'\}\mid Q)$, which is not a sum
of charges on $p$ and $p'$ that telescope separately. This limitation of Lemma~E is not an artefact: the half bound itself
fails at degree four, on a triangle-free graph of odd girth five (\cref{ex:deg4}), while \cref{thm:circular} covers graphs
of circular chromatic number below three at every degree. The
fourth paper shows that the local linear-algebraic necessary conditions on a smallest counterexample, checked inside
tree-like regions, cannot give rate $\frac12+o(1)$ at degree four or more~\cite[Proposition~6.28]{FourthPaper2026}.
\end{remark}

\begin{example}[The half bound fails at degree four]\label{ex:deg4}
\formal{DegreeFour.lean}{g13, g13_triangleFree, g13_five_cycle, g13_card_nbhd, g13code_isStorageCode, g13code_card,
g13_not_rateBound}
Let $G_{13}$ be the graph on $\{0,\dots,12\}$ in which the neighbours of $0,1,\dots,12$ are, in this order,
$\{4,7,9,10\}$, $\{5,7,8,10\}$, $\{6,7,8,9\}$, $\{8,9,10,11\}$, $\{0,8,11,12\}$, $\{1,9,11,12\}$, $\{2,10,11,12\}$,
$\{0,1,2,12\}$, $\{1,2,3,4\}$, $\{0,2,3,5\}$, $\{0,1,3,6\}$, $\{3,4,5,6\}$ and $\{4,5,6,7\}$. It is $4$-regular and
triangle-free, and $0,4,8,1,7$ is a $5$-cycle, so its odd girth is five. Over the field $\F_4=\{0,1,\omega,\omega^2\}$,
$\omega^2=\omega+1$, let $C\subseteq\F_4^{13}$ be the row space of the $7\times13$ matrix $[I_7\mid D]$, whose
column~$v$ belongs to vertex~$v$, where the rows of~$D$ are
\[
  (1,0,1,1,0,1),\ (1,1,0,1,0,1),\ (1,1,1,0,0,1),\ (0,1,1,1,1,0),\ (1,1,0,0,\omega,\omega),\ (1,0,1,0,\omega,\omega),\
  (1,0,0,1,\omega,\omega).
\]
Every column of $[I_7\mid D]$ is an $\F_4$-linear combination of the columns of the four neighbours of its vertex, so
every symbol of a word of~$C$ is a function of the symbols at its neighbours: $C$ is a storage code for $N_{G_{13}}$ over the
alphabet~$\F_4$, with $\card C=4^7$. Hence $\card C^2=4^{14}>4^{13}=\card\F_4^{\card V}$, and by \cref{thm:codes} the relation
$N_{G_{13}}$ does not satisfy \eqref{eq:star}.
\end{example}

\begin{outline}
The formal proof writes a symbol $a+b\omega$ of $\F_4$ as the bit pair $(a,b)$, so that $C$ is the image of a
$\Z/2$-linear map on $(\Z/2)^{14}$ (\lean{g13enc}, \lean{g13code}); the generator and the recovery coefficients, $2\times2$
blocks for multiplication by elements of $\F_4$, are produced by \leanfile{tools/degree_four.py}. The kernel checks that
each of the fourteen basis words satisfies the recovery identity at every vertex (\lean{g13gen_rec}); linearity extends it
to every word (\lean{g13enc_rec}), which is the storage property (\lean{g13code_isStorageCode}). The first seven symbols of
a word are its message (\lean{g13gen_id}), so the encoding is injective and $\card C=2^{14}$ (\lean{g13code_card}). With
\cref{thm:codes} this contradicts the half-rate bound (\lean{g13_not_rateBound}). Triangle-freeness, the $5$-cycle and the
degrees are decided by the kernel (\lean{g13_triangleFree}, \lean{g13_five_cycle}, \lean{g13_card_nbhd}). The alphabet
of the formal statement is $(\Z/2)^2$, of size four, as for~$\F_4$.
\end{outline}

\section{What this means for transposition}\label{sec:construction}

Our aim is to decide~(T), and we would prefer a construction. This section says what the results above require of a
linear-time transposer, and of an oblivious one in particular. Apart from \cref{prop:dn3-transfer}, a conditional
statement with a complete proof, it contains no new theorems: every claim is one of the results above, a result of an
earlier paper of the series, or explicitly labelled as a heuristic.

\subsection{The interface}\label{sec:con-interface}

For oblivious machines, the fourth paper's Theorem~D and its Corollary~4.11~\cite[Theorem~4.10 and
Corollary~4.11]{FourthPaper2026} reduce $\neg$(T) to two statements: \emph{depth reduction}, that the time--space dependency graph of every oblivious linear-time run can be cut
to bounded depth by a genie revealing $\eps L$ bits; and \emph{the crux}, that index codes for the side information of
the remaining bounded-depth cones need more than $\eps L$ bits~\cite[Definition~4.8]{FourthPaper2026}. Every fibre of an
index code is a storage code for the cones' recovery relation $N_I$~\cite[Remark~4.9(b)]{FourthPaper2026}. Hence a correct
oblivious transposer whose genie reveals $\eps L$ bits yields a binary storage code of rate at least $1-\eps$ on $N_I$;
conversely, rate bounds $\frac12+o(1)$ for all codes on such relations, together with depth reduction at some cost
$\eps<\frac12$, would rule out oblivious linear-time transposers. Depth reduction is proved for home-based work
heads~\cite{FourthPaper2026} and, in the companion paper published with this one~\cite{EighthPaper2026}, for every family
of head walks of bounded nesting load, which covers every machine whose work heads enter each cell boundedly often; it is
open in general.

\emph{Directed Recovery Structures}~\cite{DirectedRecovery2026} describes the relations $N_I$: for fixed-depth cones of
linear-time schedules, after deleting a small fraction of coordinates of high reuse (its Corollary~3.8), only
$O(L/\log d)$ coordinates $p$ have $\dist_E(L_p,R_p)\le r$, for every fixed $r$ (its Theorem~3.3 and Corollary~3.7).
Deleting these coordinates leaves a relation with $b_E(r)=0$, since deleting coordinates does not shorten paired paths;
for $r=3$ it is loopless and TT-free (\cref{lem:tt-paired}). Its Question~9.1 asks whether all codes on bounded-degree
relations with $b_E(r)=o(\card V)$ for every fixed $r$ have rate at most $\frac12+o(1)$; by its Proposition~9.2, an
affirmative answer gives this bound on all actual cone relations.

Our theorems are exact half bounds, at radius three, on classes of TT-free relations, several of them with no degree
bound. With \cref{cor:paid} they answer Question~9.1 affirmatively for every sequence of relations that, after deleting
$o(\card V)$ coordinates, lies in one of our classes. They do not decide (T): we do not know that cone relations lie in
these classes, and depth reduction is open. The degree-three case of Question~9.1 would follow from a single
combinatorial statement.

\begin{proposition}[What the directed degree-three conjecture would give]\label{prop:dn3-transfer}
\notformalized
Suppose that every loopless TT-free relation in which every vertex has at most three parents satisfies
\eqref{eq:star}; \cref{conj:P} below would imply this, by Theorem~D. Then every storage code on a relation $N$ in which
every vertex has at most three parents has rate at most $\frac12+b_E(3)/(2\card V)$. In particular,
\cite[Question~9.1]{DirectedRecovery2026} would have an affirmative answer for relations with recovery sets of size at
most three, with no bound on out-degrees and with the paired-distance condition used only at radius three.
\end{proposition}

\begin{informalproof}
Let $Z$ be the set of coordinates $p$ with $\dist_E(L_p,R_p)\le3$, so $\card Z=b_E(3)$. It contains every vertex with a
loop and the middle vertex of every transitive triangle (\cref{lem:tt-paired}; the path $L_p,R_z,L_x,R_p$ given by a
transitive triangle $x\to p\to z$, $x\to z$ needs no looplessness). A transitive triangle of $N-Z$ is one of $N$, whose
middle vertex lies in $Z$; so $N-Z$ is loopless and TT-free, and its vertices still have at most three parents. By
hypothesis it satisfies \eqref{eq:star}, and \cref{cor:paid} gives $\log_{\card\alpha}\card C\le(\card V+b_E(3))/2$.
\end{informalproof}

\subsection{Necessary conditions for a construction}\label{sec:con-necessary}

Consider an oblivious linear-time transposer and a genie that reveals $\eps L$ bits and leaves cones of bounded depth,
with recovery relation $N_I$ on the $L$ input coordinates (for home-based work heads such genies exist for every
$\eps>0$~\cite[Theorem~4.20 and Corollary~4.23]{FourthPaper2026}). The machine yields a binary storage code of rate at least $1-\eps$ on
$N_I$. In addition to the requirements of the earlier papers of the series~\cite{FirstPaper2026,SecondPaper2026,FourthPaper2026,DirectedRecovery2026},
it must satisfy the following.
\begin{enumerate}[label=(N\arabic*),leftmargin=3em]
\item\label{it:N1} \emph{No large induced sub-relation in our classes.} If $W$ is a set of coordinates such that
$N_I[W]$ satisfies \eqref{eq:star}, then $\card W\le2\eps L$. Indeed, \cref{cor:paid} with $Z=V\setminus W$ bounds the
code by $2^{L-\card W/2}$, while it has at least $2^{(1-\eps)L}$ words. This applies to every $W$ on which $N_I[W]$ is
$2$-dicolourable, or lies in one of the classes of Theorems~U3, T, PC, Corollaries~$\Tint$ and~CC, or, when symmetric, of
Theorems~N3 and~\ref{thm:circular}. For $2$-dicolourable $W$ this is a form of the classical bound of index coding by
acyclic induced sub-relations~\cite{BBJK2011}; the other classes are not covered by the acyclic bound.
\item\label{it:N2} \emph{Where a gain must live.} For every fixed $\delta>0$ there is a set $Z$ of at most
$\delta L+O(L/\log d)$ coordinates such that $N_I-Z$ is loopless and TT-free, with bounded degrees: the coordinates of
high reuse and those with paired distance at most three~\cite[Theorem~3.3 and Corollaries~3.7 and~3.8]{DirectedRecovery2026}.
By~\ref{it:N1}, every induced sub-relation of $N_I-Z$ on more than $2\eps L$ coordinates violates the hypotheses of each
theorem of \cref{sec:directed}. So it must contain vertices with at least three parents inside their strongly connected
component that lie on directed cycles among themselves (Theorem~T and Corollary~$\Tint$); it must have underlying degree
at least four somewhere (Theorem~U3); and it admits no partition into blocks as in Theorem~PC, in particular no
elimination certificate of the whole sub-relation (with the digon components as blocks, this requires a vertex with two
one-way parents, a vertex with two one-way arcs into one component, or a component of none of the three kinds). The
relation $\Ex$ of \cref{rem:X15} has the first two features but not the third: it has a global certificate.
\item\label{it:N3} \emph{Directed triangles do not help.} Native cone relations can contain linearly many disjoint directed
$3$-cycles~\cite[Proposition~3.10]{DirectedRecovery2026}, which defeat symmetric triangle-free hypotheses. Our directed
theorems allow them, so they are not by themselves a source of coding gain.
\item\label{it:N4} \emph{Gains at short odd distance do not transfer.} The Clebsch graph and the folded cubes carry codes
above rate $\frac12$, and their odd girth does not exceed their degree (\cref{rem:foldcube};
\cite[Corollary~6.11]{FourthPaper2026}). After paid degree trimming, cone relations have only $O(L/\log d)$ coordinates with
paired distance at most $r$, for every fixed~$r$~\cite[Theorem~3.3]{DirectedRecovery2026}, so such local structure can occupy only a vanishing fraction of a cone relation (heuristic
reading: a gain must come from structure at growing paired distance, which is the open territory of
Question~9.1).
\end{enumerate}

\subsection{Mechanisms that fail}\label{sec:con-fail}

The following classes of side information admit no storage code of rate above $\frac12$, for any alphabet and any
recovery functions; a construction cannot draw its coding gain from them, on more than a $2\eps$ fraction of the
coordinates.
\begin{itemize}[leftmargin=2em]
\item \emph{Cycles through vertices with two internal inputs.} Relations whose strongly connected components have no
transitive triangle and whose vertices with three or more internal parents induce an acyclic sub-relation, with any number
of parents in earlier components and any out-degree (Corollary~$\Tint$). This contains the TT-free case of the two-input
theorem of \emph{Directed Recovery Structures}~\cite[Theorem~5.1]{DirectedRecovery2026}.
\item \emph{Exchange blocks with private links.} Blocks of two-way exchanges, linked by one-way arcs that are private in the
sense of~\ref{it:P1} and~\ref{it:P2}, with blocks that are chordless subcubic, chorded paths or certified (Theorem~PC);
in particular every translation-invariant pattern $\Cay(G,\{s,s^{-1},t\})$ without transitive triangles, on every finite
group, when $\{s,s^{-1},t\}$ has three elements and does not contain~$1$ (Corollary~CC).
\item \emph{Subcubic directed structure.} Relations without transitive triangles whose underlying graph has maximum degree
three, directed triangles included (Theorem~U3).
\item \emph{Sparse or low-chromatic symmetric structure.} Symmetric side information of circular chromatic number below
three, at any degree (\cref{thm:circular}), and triangle-free subcubic symmetric side information (Theorem~N3), where
the optimal codes are fractional pairings and index codes need $\card V-\nustar$ symbols (\cref{thm:gn,thm:index}).
\end{itemize}
For impossibility proofs, \cref{sec:limits} records what fails as a method: private pieces (\cref{prop:no-private}),
minimum nonempty cores (\cref{rem:X15}), sweeps (\cref{rem:sweeps}) and global certificates alone (\cref{rem:petersen}).
Whether nested decompositions always exist at in-degree three is \cref{conj:P}.

\section{Open problems}\label{sec:open}

\subsection{Where the coding side stands}

For oblivious machines, the series reduces $\neg$(T) to depth reduction and an all-code storage bound on the cone
relations (\cref{sec:con-interface}). The table records the all-code half bounds now available at the exact radius three,
that is, on loopless TT-free relations.
\begin{center}
\small
\begin{tabular}{@{}>{\raggedright\arraybackslash}p{0.30\textwidth}>{\raggedright\arraybackslash}p{0.30\textwidth}>{\raggedright\arraybackslash}p{0.33\textwidth}@{}}
\toprule
\textbf{class} & \textbf{result} & \textbf{open nearby} \\
\midrule
symmetric, $\chi_c<3$, any degree & circular memo, \cref{thm:circular} & symmetric, $\chi_c=3$, degree $\ge4$ \\
\addlinespace
symmetric, triangle-free, degree $\le3$ & Theorem~N3; $\gn=\nustar$ (\cref{thm:gn}) & symmetric, triangle-free, degree $4$ \\
\addlinespace
underlying degree $\le3$ & Theorem~U3 & the all-code form of Theorem~S of~\cite{FourthPaper2026} \\
\addlinespace
vertices with $\ge3$ internal parents acyclic & Theorem~T, Corollary~$\Tint$ & in-degree $\le3$ in general (\cref{conj:P}) \\
\addlinespace
blocks with private links & Theorem~PC, Corollary~CC & chorded digon cycles (\cref{conj:K}) \\
\bottomrule
\end{tabular}
\end{center}
By \cref{rem:foldcube}, an exact half bound that assumes only $b_E(r)=0$ can hold at most up to maximum degree $r+1$;
at radius three, TT-freeness alone can give exact bounds at most up to degree four, and further classes need further
structure, as in Theorems~T and~PC. The transposition question needs the asymptotic form,
\cite[Question~9.1]{DirectedRecovery2026}.

\subsection{Open problems}

\begin{enumerate}[label=\arabic*.,leftmargin=2em]
\item \emph{Decide (T).} Construct a machine transposing $d\times d$ matrices with $b(d)$-bit entries in time $O(L)$, or
prove that none exists. The first paper's frozen checker accepts either answer~\cite{FirstPaper2026}.

\item \emph{The directed degree-three case.} The natural next class is all of in-degree three.

\begin{conjecture}[Conjecture P]\label{conj:P}
Every loopless TT-free relation in which every vertex has at most three parents has a decomposition
(\cref{def:piece}).
\end{conjecture}

By Theorem~D, \cref{conj:P} implies that all such relations satisfy \eqref{eq:star}, and by \cref{prop:dn3-transfer}
that the paired-distance question has an affirmative answer for recovery sets of size at most three, at any out-degree.
Every relation treated in this paper has a decomposition, including $\Ex$ (\cref{rem:X15}), and \cref{rem:X15} shows that
a proof cannot always start from a minimum nonempty core. One intermediate target is the all-code version of the fourth paper's
Theorem~S on its whole class: TT-free relations with at most three parents per vertex, in which every vertex with three
parents has no child that is not a parent. It lies between Theorem~U3 and \cref{conj:P}.

\item \emph{Chorded digon cycles.}

\begin{conjecture}[Conjecture K]\label{conj:K}
Let $K$ be a chorded digon cycle, with at most one chord into and at most one chord out of every vertex, such that $N[K]$
is TT-free. Then $N[K]$ has an elimination certificate.
\end{conjecture}

With Theorem~PC, \cref{conj:K} would give \eqref{eq:star} for every TT-free relation whose digon components are paths or
cycles and which satisfies (H1) and (H2) of \cref{cor:PC-digon}; this includes every TT-free relation in which every
vertex has exactly two digon partners, one one-way parent and one one-way child. \Cref{prop:CC} with the mirror sweep
proves it when all chords have one offset $k$ with $2k\not\equiv\pm1\pmod{\card K}$, in particular for TT-free
circulants; sweeps do not prove it in general (\cref{rem:sweeps}), but every cycle of \cref{rem:sweeps} has a
certificate with three records.

\item \emph{Degree four.} Does every triangle-free graph of maximum degree four satisfy \eqref{eq:star}, or at least have
guessing number $\card V/2+o(\card V)$ along families with growing girth? At degree five the exact bound fails (the Clebsch
graph), and \cref{rem:degree4} explains why Lemma~E alone does not reach degree four. The circular memo's theorem gives
the answer for $\chi_c<3$ at every degree.

\item \emph{The all-code paired-distance question}~\cite[Question~9.1]{DirectedRecovery2026}. Do bounded-degree
recovery relations with $b_E(r)=o(\card V)$ for every fixed $r$ admit only storage codes of rate at most
$\frac12+o(1)$, over every alphabet? With depth reduction for a class of oblivious schedules, an affirmative answer would
exclude oblivious linear-time transposers in that class (\cref{sec:con-interface}); a counterexample would not by itself
construct a transposer~\cite{DirectedRecovery2026}. Our results answer it for relations that are close, in the sense of
\cref{cor:paid}, to one of the classes above.

\item \emph{Formalization.} Formalize the topological order of strongly connected components used in
Corollary~$\Tint$, the digon-component form of Theorem~PC (\cref{cor:PC-digon}), \cref{lem:tt-paired},
\cref{cor:paid}, \cref{prop:dn3-transfer} and \cref{prop:no-private}, and turn the computations of \cref{sec:limits}
into Lean certificates.
\end{enumerate}


\appendix
\section{Index of Lean declarations}\label{app:index}

Every declaration of this paper's library cited in the text, grouped by the Lean module that contains it, relative to the namespace
\texttt{HalfRate}. All of them depend only on the axioms \texttt{propext}, \texttt{Classical.choice} and
\texttt{Quot.sound} (\cref{app:verify}). The declarations of the earlier libraries of the series that are named in
the text, in the namespaces \texttt{Transpose} of the first paper, \texttt{MotionDepthCodes} of the fourth paper and
\texttt{CircularMemo} of its circular memo, are defined and verified
there~\cite{FirstPaper2026,FourthPaper2026,CircularMemo2026}.

\paragraph{\texttt{HalfRate.Barrier}}
{\small\raggedright \texttt{barrier\_ineq}, \texttt{isolatedOut}, \texttt{log\_card\_le\_barrier}.\par}

\paragraph{\texttt{HalfRate.Blocks}}
{\small\raggedright \texttt{ChordlessSubcubic}, \texttt{ElimOn}, \texttt{ElimOn.piece}, \texttt{chordless\_block\_decomp}, \texttt{localGraph}.\par}

\paragraph{\texttt{HalfRate.Cayley}}
{\small\raggedright \texttt{card\_sq\_le\_cayley}, \texttt{cayN}, \texttt{coset\_steps}, \texttt{mem\_coset\_iff}, \texttt{rateBound\_cayley}.\par}

\paragraph{\texttt{HalfRate.Circulant}}
{\small\raggedright \texttt{circB}, \texttt{circCol}, \texttt{circulant\_sweep}, \texttt{elimOn\_circulant}, \texttt{elimOn\_cycle}, \texttt{elimOn\_ofFn}.\par}

\paragraph{\texttt{HalfRate.Corollaries}}
{\small\raggedright \texttt{ElimCert}, \texttt{ElimCert.pullback}, \texttt{ElimCert.rateBound}, \texttt{IsAcyclicIn}, \texttt{andrasfaiCert}, \texttt{rateBound\_of\_circularHom'}, \texttt{rateBound\_of\_twoAcyclic}.\par}

\paragraph{\texttt{HalfRate.Decomposition}}
{\small\raggedright \texttt{Decomp}, \texttt{PieceCert}, \texttt{PieceCert.ineq}, \texttt{RelCharged}, \texttt{RelCharged.charged}, \texttt{card\_sq\_le\_of\_decomp}, \texttt{contr}, \texttt{contr\_insert\_eq}, \texttt{contr\_isPolymatroid}, \texttt{decomp\_ineq}, \texttt{rateBound\_of\_decomp}.\par}

\paragraph{\texttt{HalfRate.IndexCoding}}
{\small\raggedright \texttt{IndexCode}, \texttt{IndexCode.isStorageCode\_fibre}, \texttt{exists\_indexCode}, \texttt{log\_card\_message\_ge}.\par}

\paragraph{\texttt{HalfRate.Matching}}
{\small\raggedright \texttt{IsFracMatching}, \texttt{dcFrac}, \texttt{dcFrac\_isFracMatching}, \texttt{exists\_code\_fracMatchingNumber}, \texttt{exists\_doubleCover\_matching}, \texttt{fracMatchingNumber}, \texttt{fracMatchingNumber\_eq\_min\_barrier}, \texttt{fracMatchingNumber\_spec}, \texttt{fracWeight}, \texttt{fracWeight\_dcFrac}, \texttt{fracWeight\_le\_barrier}, \texttt{guessingNumber\_eq\_fracMatching}, \texttt{log\_card\_le\_fracMatchingNumber}.\par}

\paragraph{\texttt{HalfRate.N3}}
{\small\raggedright \texttt{CycIn}, \texttt{CycIn.chordless}, \texttt{CycIn.share}, \texttt{block\_decomp\_of\_cycles}, \texttt{card\_sq\_le\_n3}, \texttt{cyclePiece}, \texttt{decomp\_of\_cycles}, \texttt{n3\_decomp}, \texttt{rateBound\_n3}.\par}

\paragraph{\texttt{HalfRate.NestedElimination}}
{\small\raggedright \texttt{Charged}, \texttt{defect}, \texttt{defect\_eq\_mi}, \texttt{elim\_defect}, \texttt{mi\_mono}, \texttt{nested\_elimination}, \texttt{record\_increment}, \texttt{sum\_records\_step}, \texttt{totalDefect}, \texttt{totalDefect\_le}.\par}

\paragraph{\texttt{HalfRate.PathCert}}
{\small\raggedright \texttt{ChordedPath}, \texttt{ChordedPath.elimOn}, \texttt{DownChord}, \texttt{elimOn\_path}, \texttt{pathRoot}, \texttt{pathRoot\_succ\_ne}, \texttt{rateBound\_pc}.\par}

\paragraph{\texttt{HalfRate.Pieces}}
{\small\raggedright \texttt{PieceCert.mono}, \texttt{RelCharged.mono}, \texttt{bdirPiece}, \texttt{exists\_relCharged}, \texttt{pieceCert\_of\_elim}.\par}

\paragraph{\texttt{HalfRate.Positional}}
{\small\raggedright \texttt{relCharged\_of\_cosupply}, \texttt{rev}.\par}

\paragraph{\texttt{HalfRate.TheoremT}}
{\small\raggedright \texttt{BlockOrder}, \texttt{DCyc}, \texttt{DCyc.acyclic\_of\_missing}, \texttt{DCyc.chordless}, \texttt{TTFreeOn}, \texttt{block\_decomp}, \texttt{blocks\_decomp}, \texttt{card\_sq\_le\_thmT}, \texttt{exists\_hasDCyc\_of\_sourceless}, \texttt{i5Piece}, \texttt{rateBound\_blocks}, \texttt{rateBound\_thmT}, \texttt{singlePiece}.\par}

\paragraph{\texttt{HalfRate.U3}}
{\small\raggedright \texttt{card\_sq\_le\_u3}, \texttt{dtriPiece}, \texttt{dtri\_orient}, \texttt{rateBound\_u3}, \texttt{sub\_nbhd\_und}, \texttt{triPiece}, \texttt{und}.\par}

\paragraph{Also named in the text}
{\small\raggedright \texttt{TTFree}, \texttt{mi}.\par}

\section{Building and verifying the formalization}\label{app:verify}

The Lean project is the directory \leanfile{lean/} next to this paper. It uses Lean and Mathlib, both at version v4.34.1
(\leanfile{lean-toolchain}, \leanfile{lakefile.toml}, \leanfile{lake-manifest.json}). It requires, through relative paths,
the Lean library of the fourth paper~\cite{FourthPaper2026}, in the directory
\leanfile{Curved-Motion-Depth-Reduction-and-Storage-Codes-for-Matrix-Transposition-on-Tapes/lean}, and that of its circular
memo~\cite{CircularMemo2026}, in the subdirectory \leanfile{memo/lean} of the same paper. The fourth paper's library in turn
imports the second paper's library, which imports that of the first~\cite{FirstPaper2026,SecondPaper2026}, and the
repository's shared library \texttt{lean-machine}. Keep the repository's directory layout, so that these sibling directories
stay in place.

\paragraph{Build.} With \texttt{elan} installed, run in \leanfile{lean/}:
\begin{center}
\texttt{lake exe cache get \&\& lake build}
\end{center}
The first command downloads the compiled Mathlib; the second compiles the library and the parts of the earlier libraries
that it uses. Before the first build, fetch the second paper's pinned external dependency as its README describes.

\paragraph{Axioms.} The file \leanfile{PaperAxioms.lean}, generated from the \LaTeX{} sources of this paper by
\leanfile{tools/lean_index.py}, runs \texttt{\#print axioms} on every declaration of this library named here
(\cref{app:index}):
\begin{center}
\texttt{lake env lean PaperAxioms.lean}
\end{center}
Every declaration depends at most on \texttt{propext}, \texttt{Classical.choice} and \texttt{Quot.sound}. No Lean file of this
library uses \texttt{sorry}, \texttt{admit} or \texttt{native\_decide}, or declares an \texttt{axiom}. Running
\texttt{python3 tools/lean\_index.py --index-only} in the paper's directory regenerates only the index of
\cref{app:index}.

\paragraph{Where things are.} \leanfile{NestedElimination.lean} (Lemma~E), \leanfile{Decomposition.lean} (contraction,
piece certificates and Theorem~D), \leanfile{Positional.lean} (co-supply), \leanfile{Pieces.lean} (fewer arcs, greedy
certificates, Lemma~$\Bdir$), \leanfile{Corollaries.lean} (global certificates, the pull-back, $2$-dicolourable relations and
circular chromatic number below three), \leanfile{N3.lean} (Theorem~N3 and the shortest-cycle decomposition of blocks),
\leanfile{U3.lean} (Theorem~U3), \leanfile{TheoremT.lean} (Theorem~T and the block form of Corollary~$\Tint$),
\leanfile{Blocks.lean} (internal certificates and chordless subcubic blocks), \leanfile{PathCert.lean} (Lemma~S and
Theorem~PC), \leanfile{Circulant.lean} (positional certificates and Proposition~CC), \leanfile{Cayley.lean}
(Corollary~CC), \leanfile{Barrier.lean} (the barrier bound), \leanfile{Matching.lean} (fractional matchings, the fractional
Tutte--Berge formula and the guessing number), \leanfile{IndexCoding.lean} (index codes and the broadcast rate), and
\leanfile{Examples.lean} (small sanity checks).

\paragraph{Computations.} The numerical remarks of \cref{sec:limits} (\cref{rem:X15,rem:sweeps,rem:petersen}) and the
numerical check of the folded-cube formula in \cref{rem:foldcube} are reproduced by \leanfile{tools/verify_remarks.py}
(plain Python~3, run from the paper's directory as \texttt{python3 tools/verify\_remarks.py}), described in
\leanfile{tools/README.md}; the script also checks \cref{prop:no-private} by exhaustion, and \cref{ex:D9} and the global
certificate before it step by step. They are not used in any proof, and the Lean library does not depend on them.


\bibliographystyle{plainurl}
\bibliography{references}

\end{document}
